\documentclass[a4paper,twoside,12pt]{report}
% Richard Klein (2020,2021)

% Include Packages
%\usepackage[a4paper,inner=3.5cm,outer=2.5cm,top=2.5cm,bottom=2.5cm]{geometry}  % Set page margins
\usepackage{tikz}
\usepackage[font=small]{caption}
\usepackage{subcaption}
\usepackage{fullpage}
\usepackage{float}                  % Allows 'Here and Only Here' [H] for Floats
\usepackage{url}                    % \url{} command
\usepackage{charter}                  % Set font to Times
\usepackage{graphicx}               % \includegraphics
\usepackage{amsmath}
\usepackage{amssymb}
\usepackage{amsthm}
\usepackage{booktabs}
\usepackage{parskip}
\usepackage{ragged2e}
\usepackage[all]{nowidow}
\setnoclub[2]
\setnowidow[2]

% Referencing
% Provides \Vref and \vref to indicate where a reference is.
\usepackage{varioref} 
% Hyperlinks references
\usepackage[bookmarks=true,bookmarksopen=true]{hyperref} 
% Provides \Cref, \cref, \Vref, \vref to include the type of reference: fig/eqn/tbl
\usepackage{cleveref} 
% Setup Hyperref
\hypersetup{
  colorlinks   = true,              %Colours links instead of ugly boxes
  urlcolor     = blue,              %Colour for external hyperlinks
  linkcolor    = blue,              %Colour of internal links
  citecolor    = blue                %Colour of citations
}
% Names for Clever Ref
\crefname{table}{table}{tables}
\Crefname{table}{Table}{Tables}
\crefname{figure}{figure}{figures}
\Crefname{figure}{Figure}{Figures}
\crefname{equation}{equation}{equations}
\Crefname{equation}{Equation}{Equations}

% Wits Citation Style
\usepackage{utils/natbib} \input{utils/natbib-add}
\bibliographystyle{utils/named-wits}
\bibpunct{[}{]}{;}{a}{}{}  % to get correct punctuation for bibliography
\setlength{\skip\footins}{1.5cm}
\newcommand{\citets}[1]{\citeauthor{#1}'s \citeyearpar{#1}}
\renewcommand\bibname{References}  

\pagestyle{headings}

\pagestyle{plain}
\pagenumbering{roman}

\renewenvironment{abstract}{\ \vfill\begin{center}\textbf{Abstract}\end{center}\addcontentsline{toc}{section}{Abstract}}{\vfill\vfill\newpage}
\newenvironment{declaration}{\ \vfill\begin{center}\textbf{Declaration}\end{center}\addcontentsline{toc}{section}{Declaration}}{\vfill\vfill\newpage}
\newenvironment{acknowledgements}{\ \vfill\begin{center}\textbf{Acknowledgements}\end{center}\addcontentsline{toc}{section}{Acknowledgements}}{\vfill\vfill\newpage}

\begin{document}
\RaggedRight % Enforce Wits readability rule (avoid full justification)
\pagenumbering{roman} % Roman numerals for preliminaries
\thispagestyle{empty}

% -------------------------
% TITLE PAGE
% -------------------------
\begin{center}
  % \vspace*{2cm}
  {
  \huge \bf \textsc{CRISP-NAM for Survival Analysis under Competing Risks using Risk Set Sampling}\\
  \vspace{1.5cm}
  \large \textbf{Michael Falk}\\
  3052851\\[1.5cm]
  \includegraphics[width=2cm]{utils/images/wits.jpg} % Adjust path as needed
  \vspace{1.5cm}\\
  \large A research report submitted to the Faculty of Science, University of the Witwatersrand, Johannesburg, in partial fulfilment of the requirements for the degree of Master of Science.\\[1.5cm]
  Supervised by Dr. Justine Nasejje\\[5pt]
  Co-Supervised by Dr. Steven James\\[1.5cm]
  Johannesburg, \today\\
  \vspace{0.5cm}
  \small{Ethics Clearance Number: WCSAM-2026-25}\\
  }
\end{center}
\vfill
\newpage

% -------------------------
% DECLARATION (Wits Standard Wording)
% -------------------------
\begin{declaration}
I declare that this research report is my own, unaided work. It is being submitted for the Degree of Master of Science at the University of the Witwatersrand, Johannesburg. It has not been submitted before for any degree or examination at any other University.

\vspace{2cm}
\noindent\rule{6cm}{0.4pt}\\
(Signature of candidate)\\

\vspace{1cm}
\noindent \rule{1cm}{0.4pt} day of \rule{3cm}{0.4pt} 20\rule{1cm}{0.4pt} at \rule{4cm}{0.4pt}
\end{declaration}

% -------------------------
% ABSTRACT
% -------------------------
\begin{abstract}
Deep survival models for competing risks face severe computational bottlenecks due to exhaustive risk-set evaluations. To resolve this, this study integrates a novel dynamic, entropy-guided risk-set sampling strategy into an interpretable Neural Additive Model (CRISPNAM). Evaluated across synthetic and clinical datasets (PBC, Framingham), traditional sampling methods like Nested Case-Control and Static Case-Cohort have novel performance gains along with some calibration failures under high compression. Conversely, the Entropy sampler intelligently targets high-variance subjects. By leveraging stabilized Inverse Probability Weighting, the optimal entropy configuration acts as a data-level regularizer. Regarding feature recovery, CRISPNAM successfully tracked the true linear trajectories of primary synthetic predictors (such as $x_7$ and $x_8$), though L2 regularization and SiLU activations prevented complete suppression of null-effect features (such as $x_3$ and $x_4$). In the clinical PBC cohort, the model appropriately isolated critical biomarkers like serum bilirubin (serBilir) and prothrombin time as major risk contributors, while correctly assigning negligible, flat-line impact to clinical observations like hepatomegaly and spiders. The entropy-sampled model matched the Full Cohort baseline predictive performance (achieving a C-index of 0.851 on synthetic data) while reducing computational training time by over 50\%. Ultimately, this variance driven dynamic sampling provides a computationally scalable, highly calibrated, and clinically transparent framework for large-scale competing risk analysis.
\end{abstract}

% -------------------------
% ACKNOWLEDGEMENTS
% -------------------------
\begin{acknowledgements}
I would like to acknowledge Dr. Justine Nasejje for assisting and guiding me in the creation of this study, as well as all the lecturers who imparted their knowledge and time into my understanding of the work.
\end{acknowledgements}

% -------------------------
% LISTS (TOC, Figures, Tables, Nomenclature)
% -------------------------
\tableofcontents
\newpage

% \phantomsection
% \addcontentsline{toc}{chapter}{List of Figures}
% \listoffigures
% \newpage

% \phantomsection
% \addcontentsline{toc}{chapter}{List of Tables}
% \listoftables
% \newpage

% ======================================================================
% BODY CHAPTERS
% ======================================================================
\clearpage
\pagenumbering{arabic} % Switch to Arabic numerals for main text

\chapter{Introduction}
\label{ch:introduction}

\section{Background and Motivation}

The modern world is currently navigating a highly data-driven era, where the volume, velocity, and variety of information generated across multiple domains present challenges for data management and analysis \citep{Laney2001}. From global financial transactions and sensor networks to electronic health records and longitudinal patient tracking, data are collected across a wide range of activities. While this availability of data has supported advances in deep learning \citep{LeCun2015}, it has simultaneously introduced computational challenges. Training algorithms to process and extract meaningful patterns from these datasets requires computational resources, highlighting the need for scalable processing techniques.

The computational burden is also apparent in healthcare, particularly within the domain of survival analysis. Survival analysis evaluates the time until a specific event of interest occurs, such as the onset of a disease or patient mortality \citep{Cox1972}. A fundamental component of this analysis is the concept of the \textit{risk set}, which comprises subjects who remain event-free and under observation at a given time. As the dataset grows, more individuals contribute to the risk sets, increasing the computation required to evaluate the loss function. Risk-set sampling aims to reduce this computational burden by using a subset of eligible individuals at each event time.

This challenge is further compounded in real-world clinical settings by the presence of \textit{competing risks} \citep{latouche2013competing}. In a competing risks scenario, a patient may experience an alternative, mutually exclusive event that permanently precludes the primary event of interest from occurring, for example, death from a non-cardiovascular cause before developing coronary heart disease \citep{andersen2012competing}. To model the cumulative incidence of the primary event in the presence of competing events, models can use the Fine--Gray subdistribution hazard framework \citep{FineGray1999}. This formulation retains subjects who have experienced a competing event in the modified risk set, with weights that account for right censoring \citep{FineGray1999}. As a result, a competing event does not remove an individual from subsequent risk-set comparisons. In an implementation that explicitly constructs these comparisons, retaining these contributions adds to the computational work required to evaluate the objective.

Deep survival architectures can capture non-linear relationships between covariates and event risk \citep{kvamme2019}. However, the computational cost of training these models is important during model development. Model development is an inherently iterative process: an architecture must be trained before it can be evaluated, its predictions tested, and its hyperparameters tuned. If a single training run takes an excessive amount of time, or if a model trains incorrectly and must be restarted, the entire research, evaluation, and deployment pipeline is delayed. Speeding up the training phase can therefore reduce the time required for repeated testing and validation.

To address this, the thesis applies nested case-control and static case-cohort sampling to an interpretable neural Fine--Gray model and introduces, develops, and implements a dynamic entropy-guided risk-set sampling method. By shifting away from evaluating every eligible risk-set contribution within each training mini-batch and instead training on a sampled approximation, this study investigates whether training time can be reduced while limiting the loss of predictive performance. Sampling-based approaches have previously been investigated for neural Cox models, providing motivation for examining their use in this setting \citep{kvamme2019}.

This study therefore evaluates whether dynamic sampling can provide a favourable balance between training time and predictive performance. The primary goal is to achieve faster training with only a small reduction, if any, in discrimination and predictive accuracy, while also assessing calibration and computational resource usage. By establishing a framework that samples risk-set contributions using the model's current predictions, this research aims to improve the computational efficiency of neural competing risks modelling.

\section{Research Questions and Objectives}

\subsection{Research Questions}
This study investigates the effects of risk-set sampling in Fine--Gray models for right-censored competing risks data, including a proposed entropy-guided sampling strategy. The key research questions guiding this thesis are:
\begin{itemize}
    \item How do case-cohort and nested case-control sampling designs affect discrimination, calibration, and prediction error in Fine--Gray models for competing risks compared to an unsampled mini-batch baseline?
    \item How does the proposed entropy-guided sampling strategy, which adapts its sampling probabilities using the model's current predictions, affect discrimination, calibration, and prediction error compared to the unsampled baseline, nested case-control sampling and static case-cohort sampling?
    \item How do sampling settings, including the number of controls, subcohort fraction and entropy temperature, affect training stability and predictive performance?
    \item Which sampling strategies and settings provide the greatest reduction in training time while limiting deterioration in discrimination, prediction error and calibration relative to the unsampled baseline?
    \item How do different sampling strategies affect CPU usage, memory requirements, and training duration for deep Fine-Gray models?
\end{itemize}

\subsection{Research Objectives}
To address the research questions outlined above, this study pursues the following primary objectives:
\begin{itemize}
    \item To implement interpretable deep Fine-Gray models (CRISP-NAM) with Fine-Gray censoring weights and inverse-probability sampling weights for sampled risk sets using case-cohort, nested case-control, and the proposed entropy-guided sampling strategy.
    \item To evaluate and compare the predictive performance (discrimination, calibration, and prediction error) of unsampled mini-batch, case-cohort, nested case-control, and entropy-guided sampling designs using right-censored competing risks data.
    \item To determine whether the proposed entropy-guided sampling strategy reduces training time while limiting deterioration in discrimination, prediction error and calibration relative to the unsampled baseline, and to compare this trade-off with case-cohort and nested case-control sampling.
    \item To analyze the impact of different sampling strategies on computational efficiency, including CPU usage, memory requirements, and training duration across model configurations.
    \item To analyze the impact of sampling settings, including the number of controls, subcohort fraction and entropy temperature, on training stability, C-index and Brier based prediction error.
    \item To provide practical guidance on selecting sampling strategies and settings that balance computational efficiency with predictive performance within the interpretable CRISP-NAM framework.
    \item To produce reproducible code and benchmark results to support the adoption of sampling-aware deep survival analysis methods. The complete source code and experimental pipelines are publicly hosted on GitHub \citep{falk2026crispnam}.
\end{itemize}

% \section{Thesis Organization}
% This thesis is structured to systematically introduce, develop, and evaluate the proposed sampling methodologies within the context of deep survival analysis. 
% \begin{itemize}
%     \item \textbf{Chapter 1 (Introduction):} Outlines the background, motivation, research questions, and primary objectives of the study.
%     \item \textbf{Chapter 2 (Background Theory):} Explores the theoretical foundations of competing risks, the Fine-Gray subdistribution hazard, traditional risk-set sampling strategies, and modern deep survival architectures.
%     \item \textbf{Chapter 3 (Methodology):} Details the preprocessing pipelines, the mathematical formulation of the novel dynamic entropy sampling strategy, the CRISPNAM-FG network architecture, and the quantitative evaluation metrics.
%     \item \textbf{Chapter 4 (Results and Discussion):} Presents the empirical findings across both synthetic and real wolrd datasets, highlighting the trade-offs between predictive discrimination, probability calibration, and computational efficiency. 
%     \item \textbf{Chapter 5 (Conclusion and Future Work):} Summarizes the major contributions of the study, addresses limitations, and outlines potential avenues for future research in scalable survival analysis.
% \end{itemize}
% % --------------------------------------------------------------

\section{Thesis Organization}
This thesis is structured to systematically introduce, develop, and evaluate the proposed sampling methodologies within the context of deep survival analysis. Chapter 1 outlines the background, motivation, research questions, and primary objectives of the study. Chapter 2 explores the theoretical foundations of competing risks, the Fine--Gray subdistribution hazard, traditional risk-set sampling strategies, and modern deep survival architectures. Chapter 3 details the preprocessing pipelines, the mathematical formulation of the proposed entropy-guided sampling strategy, the CRISPNAM network architecture, the sampling-aware Fine--Gray loss function, and the quantitative evaluation metrics. Chapter 4 presents the empirical findings across the synthetic, PBC and modified Framingham datasets, highlighting the trade-offs between predictive discrimination, probability calibration, and computational efficiency. Finally, Chapter 5 summarizes the major contributions of the study, addresses limitations, and outlines potential avenues for future research in scalable survival analysis.

\chapter{Background Theory}
\label{ch:background_theory}

This chapter establishes the theoretical foundations required to model time-to-event data in the presence of multiple, mutually exclusive outcomes. It first reviews the mathematical challenges introduced by competing risks and describes classical statistical frameworks, including cause-specific and subdistribution hazards. The chapter then explores recent advancements in deep survival analysis, detailing how neural architectures can be adapted for competing events. Finally, it introduces the CRISPNAM framework and the risk-set sampling strategies investigated in this study to improve computational efficiency during model optimization.

\section{Competing Risks in Survival Analysis}
While standard single-event survival analysis focuses on one event of interest, clinical settings frequently include multiple competing events \citep{andersen2012competing}. For instance, a patient may die from an unrelated cause before developing the disease of interest, thereby preventing that disease from occurring \citep{andersen2012competing, prentice1978analysis}.

In cause-specific Cox regression, individuals who experience a competing event are treated as censored at the time of that event and are removed from the target cause's risk set. This treatment is appropriate for estimating the cause-specific hazard, which describes the instantaneous event rate among individuals who remain event-free. However, estimating cumulative incidence requires accounting for the hazards of all competing causes. Treating competing events as censored and using one minus the Kaplan--Meier survival estimate instead overestimates the target event's cumulative incidence \citep{andersen2012competing}.

In contrast, the Fine--Gray model retains individuals who have experienced a competing event in its modified risk set, with weights that account for right censoring. Although these individuals can no longer experience the target event, their retained contributions allow the model to describe the subdistribution hazard associated with its cumulative incidence \citep{FineGray1999}.

\section{Past Deep Neural Network Models}
Advances in deep learning have extended survival models through neural architectures capable of representing non-linear relationships between covariates and event risk \citep{Katzman2018, kvamme2019}. A foundational model in this domain is DeepSurv, which extended the traditional Cox proportional hazards framework by replacing the standard linear predictor with a deep neural network \citep{Katzman2018}. The network learns the relationship between a patient's covariates and their risk by optimizing the Cox partial-likelihood objective with parameter regularisation \citep{Katzman2018}. Building upon this neural-Cox integration, \cite{kvamme2019} developed both proportional and non-proportional hazards models, allowing the relationship between covariates and risk to vary over time in the latter. Their approach also uses sampling derived from nested case-control methodology to approximate the Cox partial likelihood during training. These formulations address single-event survival rather than directly modelling the joint distribution of event time and competing event type \citep{Katzman2018, kvamme2019}.

To address multiple event types, methodologies such as DeepHit \cite{Lee2018} have been developed. DeepHit presents a joint modelling framework that employs a shared representation network coupled with cause-specific sub-networks to model the joint distribution of event time and type. The model optimization combines negative log-likelihood with a ranking loss to encourage concordance between predicted and observed outcomes \citep{Lee2018}. However, this architecture does not explicitly decompose predictions into separate contributions from individual input features \citep{Lee2018}. This motivates examining additive neural architectures when feature-level interpretation is an objective of the study.

\section{Cause Specific Hazard}
A classical approach to competing risks analysis is to model the cause-specific hazard for each event type \citep{prentice1978analysis}. The proportional hazards framework introduced by \cite{Cox1972} can be applied separately to each cause-specific hazard. Under this formulation, the model evaluates the instantaneous rate of the primary event, conditional on the patient having survived completely event-free up to that time.

For each cause $k$, the cause-specific hazard $\lambda_k(t|x)$ is mathematically defined as the instantaneous rate of experiencing event type $k$ at time $t$, given that the individual has survived without any event up until time $t$:
$$\lambda_k(t|x) = \lim_{\Delta t \to 0} \frac{P(t \le T < t+\Delta t, E = k \mid T \ge t, x)}{\Delta t} \,.$$

The cause-specific hazards can be used to derive the cumulative incidence function (CIF), which gives the probability that event type $k$ occurs by time $t$ in the presence of competing risks \citep{prentice1978analysis, pintilie2006competing}. It is computed by integrating the cause-specific hazard weighted by the overall event-free survival function $S(s|x)$:
$$F_k(t|x) = \int_0^t \lambda_k(s|x) S(s|x) ds \,.$$

Here, $S(s|x)$ is the probability of remaining free of all event types up to time $s$, so it depends on the hazards of all competing causes. Consequently, the cumulative incidence of event $k$ depends on both its own cause-specific hazard and those of the competing events \citep{andersen2012competing}.

When fitting a cause-specific Cox model, subjects who experience a competing event are treated as censored at its occurrence and removed from the risk set. This is appropriate because the cause-specific hazard is defined among individuals who remain event-free. However, this treatment does not mean that a competing event is equivalent to ordinary loss to follow-up: after a competing event, the target event can no longer occur \citep{andersen2012competing}.

The limitation arises when the target cause's hazard alone is used to estimate its cumulative incidence without accounting for competing causes. Cause-specific models can provide valid cumulative incidence estimates when the hazards of all causes are combined through the overall event-free survival function. Fine--Gray modelling offers an alternative by modelling the subdistribution hazard and retaining individuals with prior competing events in its modified risk set \citep{andersen2012competing, FineGray1999}.

\section{Fine-Gray Subdistribution Hazard}

The Fine--Gray subdistribution model provides an alternative to cause-specific hazard modelling by specifying a regression model directly linked to the cumulative incidence of the target event \citep{FineGray1999}. This method models the subdistribution hazard, maintaining subjects who experience competing events in the modified risk set.

As an example of this contrast, in cause-specific Cox models, competing events are treated as censored and the individuals are removed from the risk set. In the Fine--Gray model, individuals who have experienced competing events remain in the modified risk set. The subdistribution risk pool includes both subjects who are event-free at time $t$ and those who experienced a competing event before time $t$. When follow-up is subject to right censoring, the contributions of retained competing-event subjects are adjusted using censoring weights. In this study, these weights are calculated from a Kaplan--Meier estimate of the censoring survival function. They account for censoring rather than restoring the possibility of the target event occurring after a competing event \citep{FineGray1999}.

Under the proportional subdistribution hazards assumption, the Fine--Gray model is expressed as:
$$\tilde{\lambda}_k(t|X)=\tilde{\lambda}_{k,0}(t)\exp(\beta^\top X) \,,$$
where $\tilde{\lambda}_{k,0}(t)$ is the baseline subdistribution hazard, which captures how the subdistribution hazard for event $k$ changes over time, and $\exp(\beta^\top X)$ scales this baseline according to the individual's covariates $X$. The exponential term represents the individual's subdistribution hazard relative to the baseline \citep{FineGray1999}.

By retaining individuals who have experienced competing events in the modified risk set, the \cite{FineGray1999} model accounts for competing outcomes through the subdistribution hazard, which is directly linked to the cumulative incidence function (CIF). As in the cause-specific approach, the CIF represents the probability that event $k$ occurs by time $t$ in the presence of competing risks. The approaches differ in how this probability is modelled. It is mathematically expressed as:
$$F_k(t|X) = P(T \le t, \Delta=k|X) \,.$$

For continuous event times, its relationship with the subdistribution hazard is:
$$F_k(t|X) = 1-\exp\left[-\int_0^t \tilde{\lambda}_k(s|X)\,ds\right] \,.$$
Unlike the cause-specific formulation, this expression links the target event's CIF to its subdistribution hazard without requiring separate models for every competing cause \citep{FineGray1999}.

A key feature of the Fine--Gray approach is the use of a modified risk set, denoted by an indicator such as $Y_i^*(t)$. This indicator identifies eligible individuals, including those who remain under observation and event-free and those who previously experienced a competing event. Individuals previously censored or having already experienced the target event are excluded. The censoring weight $w_i(t)$ adjusts the contribution of each eligible individual \citep{FineGray1999}.

Together, the eligibility indicator, censoring weight and modelled relative hazard determine $\bar{x}(t; \beta)$, which represents the risk-set weighted average of the covariates at time $t$. The censoring weights $w_i(t)$ and relative-hazard terms $\exp(\beta^\top X_i)$ are separate factors. Their product determines each individual's contribution to this average, which appears in the score function for estimating the regression parameters \citep{FineGray1999}:
$$\bar{x}(t; \beta)=\frac{\sum_i Y_i^*(t) w_i(t) X_i \exp(\beta^\top X_i)}{\sum_i Y_i^*(t) w_i(t) \exp(\beta^\top X_i)} \,.$$

\section{CRISPNAM}

While traditional deep learning approaches offer high predictive capacity, their predictions are not always directly decomposable into individual feature contributions. A foundational work shaping the theoretical framework of this thesis is the approach by \cite{InterpretableFineGray2025}, which introduced an interpretable Fine--Gray deep survival model designed specifically for competing risks. Building upon this need for transparency in deep survival analysis, the CRISPNAM architecture uses a Neural Additive Model (NAM) framework to represent non-linear feature effects while retaining an additive predictor \citep{InterpretableFineGray2025}.

Central to the CRISPNAM architecture is the use of the FeatureNet. Unlike standard fully connected neural networks that process all inputs simultaneously through hidden layers, the additive framework assigns a dedicated, independent multi-layer perceptron, called a FeatureNet, to each individual covariate. Let $X = (x_1, x_2, \dots, x_p)$ represent the set of $p$ input features. Each feature $x_j$ is passed exclusively through its own FeatureNet, which learns a feature-specific representation \citep{InterpretableFineGray2025}. In this study's implementation, each FeatureNet produces a scalar contribution $f_j(x_j)$ for the target cause, and the model is fitted separately for each target cause considered.

The additive framework mathematically combines these learned functions by summing them together with a global bias term $b$ to construct the overall log-relative subdistribution hazard predictor $f(X)$:
$$f(X) = b + \sum_{j=1}^p f_j(x_j) \,.$$

This formulation integrates with the Fine--Gray method. By replacing the traditional linear predictor $\beta^\top X$ with the learned additive neural predictor $f(X)$, the modified subdistribution hazard becomes:
$$\tilde{\lambda}_k(t|X) = \tilde{\lambda}_{k,0}(t)\exp\left(b + \sum_{j=1}^p f_j(x_j)\right) \,.$$
Here, the predictor and its feature functions refer to the chosen target cause $k$.

Because the model computes the log-relative subdistribution hazard additively, the contribution of any single feature $x_j$ to this predictor can be isolated and visualized. This architecture allows the model to capture non-linear effects of individual covariates while showing how those effects combine in the prediction \citep{InterpretableFineGray2025}. These contributions describe the fitted model's associations, rather than establishing causal effects.

\section{Risk Set Sampling Strategies}

Evaluating every eligible observation at each target-event time can make loss computation expensive. Risk-set sampling aims to reduce this computation by selecting a subset of eligible observations. Sampling-based loss approximations have previously been investigated for neural Cox models \citep{kvamme2019}. In this thesis, the sampling strategies are applied to the modified Fine--Gray risk sets within each training mini-batch. Let $R(t_i)$ denote the eligible observations for target-event case $i$, and let $n(t_i)=|R(t_i)|$. This set includes observations still under observation immediately before $t_i$, as well as those with a prior competing event. Prior competing events remain eligible and contribute with the Fine--Gray censoring weights described earlier.

One such method is Nested Case-Control (NCC) sampling, a time-matched design where, for each observed target-event case $i$ at a given event time $t_i$, up to $m$ controls are sampled uniformly without replacement from $R(t_i)$, excluding the current case. Because controls are matched sequentially to specific events, each event has its own sampled risk set, $\tilde{R}(t_i)$. The current case is always included. The probability of a control $j$ being included depends on the number of available eligible controls. This inclusion probability is mathematically expressed as:
\[
P\!\left(j\in\tilde{R}(t_i)\mid j\in R(t_i),\,j\neq i\right)
=
\frac{\min\{m,n(t_i)-1\}}{n(t_i)-1},
\qquad n(t_i)>1.
\]

If fewer than $m$ controls are available, all available controls are included. If no controls are available, the sampled risk set contains only the current case. Sampling without replacement applies within each event-specific risk set; an observation may still be selected for multiple event times.

Alternatively, case-cohort sampling selects a random subcohort independently of event times \citep{Prentice1986}. The Static Case-Cohort (SCC) implementation in this thesis follows this principle by independently selecting each training-fold observation with probability $\alpha$, once before training. This static subcohort, denoted by $C$, is then reused throughout the entirety of the training process. At any given event time $t_i$, the evaluated risk set consists of the current case $i$ along with any members of the fixed subcohort who are eligible under the Fine--Gray risk-set definition and present in the mini-batch:
\[
\tilde{R}(t_i)=\{i\}\cup\bigl(C\cap R(t_i)\bigr).
\]

Under this design, the probability of any individual being initially selected into the subcohort is given by:
\[
P(j\in C)=\alpha.
\]

Thus, $\alpha$ determines the expected subcohort fraction rather than an exact subcohort size. Membership remains fixed across training batches and epochs. The current event case is always included, regardless of subcohort membership; other eligible observations are included only if they belong to the subcohort. This also applies to observations with prior competing events.

\section{Inverse Probability Weighting (IPW) and Stabilization}

Risk-set sampling retains each target-event case while selecting only a subset of eligible controls. To account for these unequal inclusion probabilities, Inverse Probability Weighting (IPW) is applied \citep{HorvitzThompson1952}.

This technique scales the contribution of each selected observation by the reciprocal of its sampling probability. If $\pi_i$ represents the probability that individual $i$ is selected into the sample, the standard IPW is expressed mathematically as:
$$w_i = \frac{1}{\pi_i} \,.$$
Consequently, individuals who possess a lower likelihood of being sampled are assigned proportionally larger weights. In risk-set sampling, these probabilities may depend on the event-specific risk set. This weighting accounts for sampling in the risk-set contributions, but does not by itself guarantee unbiased estimates of the model parameters or predictions.

However, small inclusion probabilities can produce large inverse probability weights and increase variability in the sampled loss contributions.

The motivation for stabilisation in this thesis was to address these contributions directly, rather than relying on L2 regularisation of the model parameters. In the current implementation, the weighted risk-set denominator is evaluated using log-sum-exp for numerical stability. This calculation does not normalise the sampling weights or remove sampling variability: the sampling correction remains the reciprocal inclusion probability.

Beyond intentional sampling disparities, survival datasets are routinely subject to right-censoring, where the exact time of the event is not recorded for all individuals. The Fine--Gray loss handles censoring separately through censoring-adjusted risk-set weights, including the weights assigned to observations with prior competing events \citep{FineGray1999}. In this implementation, the censoring survival function, $\hat G(t)$, is estimated from the training fold.

The loss combines the Fine--Gray censoring weight with the inverse sampling probability for each selected eligible observation. These components serve different purposes: one accounts for censoring in the modified risk set, while the other accounts for risk-set sampling. Their precise formulation is given in the loss-function section.



\chapter{Methodology}
\label{ch:methodology}

This chapter details the experimental design, data processing pipelines, and computational frameworks used to evaluate the proposed CRISPNAM architecture along with the  incoperation of the smapling techniques, that also include our custome sampling strategy. It begins by outlining the acquisition and preprocessing of both real wolrd clinical datasets and synthetically generated cohorts designed to test specific mathematical assumptions. The chapter then details the implementation of the deep survival models, including the network architecture, loss function formulations, and the risk set sampling techniques. Finally, it establishes the evaluation metrics and cross-validation strategies used to rigorously benchmark predictive performance .

\section{Data Preprocessing}

To evaluate the robustness, accuracy, and interpretability of the competing risks models, experiments were conducted across three distinct data environments: a fully controlled synthetic dataset, the clinical Primary Biliary Cholangitis (PBC) dataset, and a modified version of the Framingham Heart Study dataset.

\subsection{Synthetic Data Generation}
A fully synthetic dataset was computationally generated to provide a controlled environment where the true underlying data generating distributions, baseline hazards, and covariate coefficients were known definitively. The simulation is grounded in the cause-specific Cox proportional hazards framework for competing risks \citep{Cox1972, prentice1978analysis}. As theoretically formalized earlier in Chapter 2, this approach models a separate proportional hazard for each event type utilizing constant cause-specific baseline hazards. This ground-truth knowledge allows for exact empirical validation of the model's ability to recover true risk estimators without the confounding variables inherent to real wolrd medical data.

A cohort of $N = 15,000$ independent samples was generated. To ensure full experimental reproducibility, a fixed global random state (seed = 42) was utilized throughout the generation process. The feature space was constructed using 12 covariates: six continuous variables drawn from a standard normal distribution $\mathcal{N}(0, 1)$, and six binary variables sampled from a binomial distribution $\text{Binomial}(1, 0.5)$. These were concatenated to form the complete covariate matrix $X$.

Two competing events were simulated. For each event type $k \in \{1, 2\}$, a specific vector of true coefficients $\beta_k$ was predefined. To test the model's capacity for intrinsic feature selection and to simulate uninformative clinical noise, several covariates in both $\beta_1$ and $\beta_2$ were deliberately assigned a true coefficient of zero. The linear predictors for each individual were calculated as $\eta_{ik} = X_i^\top \beta_k$.

The cause-specific hazards were then defined exponentially as
$$\lambda_k(t|X_i) = \lambda_{0,k} \exp(\eta_{ik}) \,,$$
where the baseline hazards were fixed at $\lambda_{0,1} = 0.01$ and $\lambda_{0,2} = 0.05$. 

The total instantaneous hazard for any event occurring was computed as the sum of the cause-specific hazards ($\lambda_{\text{total}} = \lambda_1 + \lambda_2$). The exact time to event for each subject was simulated via inverse transform sampling, drawing from an exponential distribution parameterized by the total hazard:
$$T_{\text{event}} = \frac{-\ln(U)}{\lambda_{\text{total}}} \,,$$
where $U \sim \text{Uniform}(0,1)$. If an event occurred, the specific event type (1 or 2) was assigned probabilistically using a binomial trial, where the probability of experiencing Event 1 was strictly proportional to its relative contribution to the total hazard ($\lambda_1 / \lambda_{\text{total}}$) \citep{beyersmann2009simulating}.

Finally, right-censoring was introduced independently to mimic study dropout or administrative study end. Censoring times, $T_{\text{censor}}$, were drawn from an exponential distribution with a predefined censoring rate of $0.01$. The final observed survival time was derived as the minimum of the simulated event time and the censoring time ($\tilde{T} = \min(T_{\text{event}}, T_{\text{censor}})$), and the event indicator was recorded as 0 if the individual was censored, or assigned the respective event type otherwise.

\subsection{Primary Biliary Cholangitis (PBC) Dataset}
The Primary Biliary Cholangitis (PBC) dataset, sourced from the Mayo Clinic trial \cite{fleming1991counting, murtaugh1994primary}, provides a classical real-world baseline for evaluating competing risks. Data preprocessing for this longitudinal cohort initially required standardizing the temporal scale; observation times were converted from years into days to provide a higher-resolution continuous time axis for the neural network. 

The clinical outcomes were mapped into a defined competing risks state space: State 0 denoted right-censored patients, State 1 denoted death (the primary event of interest), and State 2 denoted liver transplantation (the competing event). Addressing missingness in the clinical covariates was crucial for tensor continuity. Missing continuous variables (such as serum cholesterol or prothrombin time) were imputed using the cohort mean, and categorical covariates were imputed using the mode (the most frequent class) \citep{little2019statistical}. These categorical variables were subsequently transformed via one-hot encoding \citep{bishop2006, kuhn2013applied}. 

Finally, while redundant administrative tracking columns and records lacking valid survival times were systematically excluded, the unique patient identifier was explicitly retained. Preserving this identifier within the final analytical cohort was essential to facilitate rigorous, patient-grouped cross-validation and prevent subject-level data leakage across repeated longitudinal visits.

\subsection{Framingham Heart Study Dataset}
To evaluate the models on a highly dimensional and well-understood cardiovascular cohort, the standard Framingham Heart Study dataset was adapted \citep{InterpretableFineGray2025}. Initial data cleaning involved resolving missing values via median imputation \cite{little2019statistical}, ensuring that feature distributions were not severely skewed by outliers. 

Because standard Framingham data is often structured for binary classification, famously popularized as predicting a binary 10-year coronary heart disease outcome rather than continuous time-to-event analysis \cite{wilson1998prediction}, the original deterministic target column (\texttt{ten\_year\_c\_h\_d}) was explicitly dropped from the feature matrix to eliminate any risk of data leakage during model optimization. As the baseline dataset lacked precise longitudinal tracking for multiple event types, time-to-event data was  generated to classify patient trajectories into three mutually exclusive states.

Event 1 (The Main Event) represented patients who actively developed Coronary Heart Disease (CHD). Their time-to-event was distributed randomly within the 10-year observational window. To induce a mathematically valid competing risk scenario, Event 2 (The Competing Event) was engineered to simulate a scenario where approximately 15\% of the non-CHD patients experienced a terminal occurrence such as a fatal accident or non-cardiovascular mortality prior to the 10-year mark. 

This  modification is vital for assessing the Fine-Gray loss function, as experiencing this alternative Event 2 permanently precludes the patient from ever experiencing the main CHD event, fundamentally shifting the risk set dynamics. Finally, Event 0 (Censoring) captured the remaining baseline patients who either safely survived event-free to the end of the 10 year study or were randomly lost to follow-up.

\section{Risk Set Sampling}

\subsection{Full Risk}
The Full Risk (FULL) baseline evaluates all eligible observations within each training mini-batch at every target-event time. The eligible set includes observations with $T_i \ge t_j$, together with observations that experienced a competing event before $t_j$. No additional risk-set sampling is performed, so each eligible observation has inclusion probability 1, while the Fine--Gray censoring weights remain part of the loss. FULL therefore provides an unsampled mini-batch baseline against which the sampling strategies are compared; it does not evaluate the entire training cohort at each optimisation step.


\subsection{Case Cohort}
The static case-cohort design offers an alternative that operates independently of event times \citep{Prentice1986}. In this architecture, a fixed subcohort is randomly selected once per training fold, before training, by independently including each observation with probability $\alpha$, and systematically reused across the entire training lifecycle. Membership remains fixed across mini-batches and epochs, with stable observation identifiers preserving membership when the data are shuffled. At each target-event time, the current case is always retained, while controls are drawn exclusively from eligible subcohort members within the current mini-batch. Eligibility includes prior competing events under the Fine--Gray risk-set definition. Other future target-event cases are not automatically included as controls unless they belong to the subcohort. The current case therefore has inclusion probability $\pi_{ii}=1$, while each other eligible observation has sampling inclusion probability $\pi_{ij}=\alpha$. Here, $\alpha$ determines the expected subcohort fraction rather than an exact subcohort size. These inclusion probabilities enter the loss through inverse sampling weights, which are combined with the separate Fine--Gray censoring weights.


\subsection{Nested Case Control}
Conversely, nested case-control sampling employs a time-dependent matching strategy. At each target-event time $t_i$, the observed case is paired with up to $m$ controls, which are drawn uniformly without replacement from the eligible Fine--Gray risk set within the current mini-batch, excluding the current case. This includes observations with prior competing events. If $N_i$ eligible controls are available, each has inclusion probability $\pi_{ij}=\min(m,N_i)/N_i$, for $N_i>0$. The current case is always retained with inclusion probability 1. If no eligible controls are available, only the current case is retained. The selected controls contribute to the loss through inverse probability weighting \citep{HorvitzThompson1952}, combined with the separate Fine--Gray censoring weights.



\subsection{Dynamic Sampling Strategy}

Moving beyond classical methodologies, recent literature has begun exploring scalable algorithms for large competing risks datasets \cite{kawaguchi2020scalable} and combining traditional sampling designs with the Fine--Gray model \cite{kim2020stratified}. Building upon these foundations, this thesis introduces a novel, adaptive, entropy-guided sampling framework applied directly to the Fine--Gray subdistribution risk set. The proposed strategy operates within each training mini-batch, where individuals who have previously experienced competing events remain eligible for selection according to the modified Fine--Gray risk-set definition.

Unlike traditional sampling methods, which select controls randomly or according to fixed sampling fractions, the proposed method evaluates the informational characteristics of individual observations using the neural network's current risk predictions.

For each individual $j$, the predicted neural log-relative subdistribution hazard score $\hat{\eta}(\boldsymbol{X}_j)$ is first transformed using the sigmoid function:

\begin{equation}
p_j =
\sigma\left(\hat{\eta}(\boldsymbol{X}_j)\right)
=
\frac{1}{1+\exp(-\hat{\eta}(\boldsymbol{X}_j))}.
\end{equation}

This transformation maps the unbounded neural risk score onto the interval $(0,1)$, allowing a bounded measure of information to be calculated. The transformed value is used as a sampling score rather than being interpreted as a calibrated probability of experiencing the event of interest.

The informational value of each individual is subsequently evaluated using normalized binary Shannon entropy:

\begin{equation}
H_j =
\frac{
-p_j\log(p_j)
-(1-p_j)\log(1-p_j)
}{\log 2}.
\end{equation}

The division by $\log 2$ normalizes the binary entropy by its theoretical maximum, ensuring that $H_j \in [0,1]$. Entropy reaches its maximum when $p_j=0.5$, representing the greatest uncertainty in the transformed prediction, and approaches zero as $p_j$ approaches either zero or one. Consequently, individuals with higher entropy scores receive greater sampling preference. Here, entropy represents a heuristic measure of uncertainty in the transformed neural risk score, rather than directly measuring clinical outcome uncertainty.

To regulate the strength of this sampling preference, an importance score is assigned to each individual using a temperature parameter $\tau>0$:

\begin{equation}
a_j =
\exp\left(\frac{H_j}{\tau}\right).
\end{equation}

The temperature parameter controls the concentration of the importance scores. Smaller values of $\tau$ amplify differences between entropy scores, placing greater emphasis on high-entropy individuals, whereas larger values reduce these differences and produce a sampling distribution closer to uniform selection.

For an observed event of interest at time $T_i$, let $C_{\mathcal{B}}(T_i)$ denote the eligible controls within the modified Fine--Gray risk set of the current mini-batch $\mathcal{B}$, excluding the event case itself. The average importance score among these controls is defined as:

\begin{equation}
\bar{a}(T_i) =
\frac{1}{|C_{\mathcal{B}}(T_i)|}
\sum_{j\in C_{\mathcal{B}}(T_i)} a_j.
\end{equation}

Rather than enforcing a fixed number of controls, the importance score of each individual is scaled relative to the average importance of the eligible controls. This produces an individual inclusion probability:

\begin{equation}
\pi_{ij} =
\min\left(
1,\;
\max\left(
\epsilon,\;
\rho
\frac{a_j}{\bar{a}(T_i)}
\right)
\right),
\end{equation}

where $\rho\in(0,1]$ represents the baseline retention rate and $\epsilon>0$ is a small numerical constant ensuring positive inclusion probabilities. The parameter $\rho$ governs the overall level of sampling compression, while $\tau$ controls the relative preference for high-entropy observations.

The equations above describe the sampling rule without numerical safeguards; the implementation also adds small positive constants inside the entropy logarithms and to the denominators used to calculate the mean importance score and inclusion probabilities.

Each eligible control is subsequently selected independently using a Bernoulli distribution:

\begin{equation}
S_{ij} \sim \operatorname{Bernoulli}(\pi_{ij}).
\end{equation}

The observed event case is always retained, with $\pi_{ii}=1$. If no eligible controls are available, only the event case is retained. The resulting sampled risk set, denoted by $\widetilde{\mathcal{R}}_{\mathcal{B}}(T_i)$, therefore contains the current event case together with the controls selected through entropy-guided Bernoulli sampling.

Through this mechanism, the number of selected controls can vary across event-specific risk sets, whereas nested case-control sampling selects $m$ controls whenever at least $m$ eligible controls are available. The number of retained risk-set observations changes according to their individual inclusion probabilities. Furthermore, because these probabilities depend on the neural network's current predictions, the sampling distribution adapts as the model parameters are updated throughout training.

The contribution of the sampled risk set to the Fine--Gray partial-likelihood denominator is estimated using inverse-probability weighting:

\begin{equation}
\widehat{D}_{\mathcal{B}}(T_i)
=
\sum_{j\in\widetilde{\mathcal{R}}_{\mathcal{B}}(T_i)}
\frac{w_{ij}^{FG}}{\pi_{ij}}
\exp\left(\eta(\boldsymbol{X}_j)\right),
\end{equation}

where $w_{ij}^{FG}$ denotes the Fine--Gray censoring adjustment and $\pi_{ij}$ represents the individual's inclusion probability. The corresponding event-specific negative log-partial-likelihood contribution is:

\begin{equation}
\widehat{\mathcal{L}}(T_i)
=
-\left[
\eta(\boldsymbol{X}_i)
-
\log\widehat{D}_{\mathcal{B}}(T_i)
\right].
\end{equation}

This formulation adjusts the sampled risk-set contributions according to their inclusion probabilities, while retaining the censoring correction required by the Fine--Gray model. The complete weighting formulation and mini-batch loss function are developed in Section~3.4.

Because the sampling probabilities are derived from the actively updating neural network, the proposed strategy functions as an adaptive risk-set approximation. Although inverse-probability weighting provides a mechanism for correcting unequal inclusion probabilities in the risk-set denominator, it does not independently guarantee unbiasedness of the resulting partial-likelihood estimator. The effectiveness of this strategy is therefore investigated empirically, evaluating its ability to maintain predictive performance while reducing computational requirements relative to the alternative sampling methods.



\subsection{IPW and Stability}

Fundamentally, all evaluated sampling strategies, including case-cohort, nested case-control, and the proposed dynamic entropy-guided design, are unified through the Inverse Probability Weighting (IPW) framework. Under each sampling design, selected individuals contribute to the Fine--Gray risk-set denominator inversely proportional to their probabilities of inclusion.

Let $\pi_i$ represent the inclusion probability of an individual $i$. Under correctly specified positive inclusion probabilities, the Horvitz--Thompson \cite{HorvitzThompson1952} principle establishes that
\begin{equation}
    \mathbb{E}\left[
        \frac{\mathbf{1}\{i \in S\}}{\pi_i}\phi_i
    \right]
    = \phi_i,
\end{equation}
where $\phi_i$ represents a fixed individual contribution to a complete risk-set quantity. This establishes the theoretical justification for applying IPW to recover additive risk-set contributions in expectation.

However, although inverse-probability weighting corrects unequal representation in expectation, small inclusion probabilities can produce large weights and increase estimator variability. To support numerical stability, the implementation enforces strictly positive inclusion probabilities, evaluates weighted risk-set denominators in the logarithmic domain using the log-sum-exp operation, and applies gradient clipping during neural-network optimization.

These numerical safeguards do not guarantee unbiasedness or consistency of the resulting nonlinear partial-likelihood estimator. Consequently, the effectiveness and stability of the weighting framework are evaluated empirically across the different sampling strategies.
\section{Network Architecture (CRISPNAM-FG)}
\subsection{Network architecture}

The modeling framework follows the CRISP-NAM design \citep{InterpretableFineGray2025}, extended here to accommodate mini-batch risk-set sampling under the competing-risks Fine Gray formulation. The architecture is grounded in a Neural Additive Model (NAM) structure, wherein each covariate acts on the sub-distribution hazard independently through its own dedicated sub-network, termed a \textit{FeatureNet}.

Let $\mathbf{x} = (x_1, x_2, \dots, x_p)^\top \in \mathbb{R}^p$ denote the vector of input covariates for a subject. Each scalar feature $x_i$ ($i = 1, \dots, p$) is mapped through a dedicated feedforward sub-network $f_i(x_i): \mathbb{R} \to \mathbb{R}$, yielding an isolated non-linear contribution to the cause-specific log-subdistribution hazard.

\subsubsection{FeatureNet Formulation}
Each FeatureNet $f_i(\cdot)$ is constructed as a deep multilayer perceptron comprising $L$ hidden layers. Setting the initial layer input to the scalar covariate $z_i^{(0)} = x_i$, the transformation proceeds recursively across hidden layers $l = 1, \dots, L$:
\begin{align}
    u_i^{(l)} &= W_i^{(l)} z_i^{(l-1)} + b_i^{(l)}, \\
    z_i^{(l)} &= \text{Dropout}\left( \text{SiLU}\left(u_i^{(l)}\right); p_{\text{drop}} \right), \quad l = 1, \dots, L,
\end{align}
where $W_i^{(l)} \in \mathbb{R}^{d_l \times d_{l-1}}$ and $b_i^{(l)} \in \mathbb{R}^{d_l}$ denote the layer weights and bias vectors, respectively, and $p_{\text{drop}}$ is the dropout retention rate. The sub-networks utilize the Sigmoid Linear Unit ($\text{SiLU}$) activation function, defined as:
\begin{equation}
    \text{SiLU}(u) = u \cdot \sigma(u) = \frac{u}{1 + e^{-u}}.
\end{equation}
The smooth, continuously differentiable, and non-monotonic nature of $\text{SiLU}$ mitigates vanishing gradients across deeper sub-networks while yielding smooth univariate feature effect curves.

The final representation $z_i^{(L)} \in \mathbb{R}^{d_L}$ is projected to a scalar output via a bias-free linear transformation:
\begin{equation}
    \tilde{f}_i(x_i) = W_i^{(L+1)} z_i^{(L)},
\end{equation}
where $W_i^{(L+1)} \in \mathbb{R}^{1 \times d_L}$. To guarantee model identifiability a core requirement of generalized additive models each FeatureNet enforces a zero-mean batch constraint during training:
\begin{equation}
    f_i(x_{b,i}) = \tilde{f}_i(x_{b,i}) - \frac{1}{B} \sum_{j=1}^B \tilde{f}_i(x_{j,i}),
\end{equation}
for sample $b$ within a mini-batch of size $B$. This centering operation ensures that individual shape functions $\mathbb{E}[f_i(x_i)] = 0$ capture purely relative variations, preventing arbitrary vertical translation between subnetworks.

\subsubsection{Global Additive Aggregation}
The overall model aggregates the individual feature-specific transformations additively with a learnable global intercept $\beta_0 \in \mathbb{R}$:
\begin{equation}
    \eta(\mathbf{x}) = \beta_0 + \sum_{i=1}^p f_i(x_i),
\end{equation}
where $\eta(\mathbf{x})$ denotes the predicted log-subdistribution hazard score (relative risk). Because the final projection layers of the FeatureNets omits additive biases, the global baseline risk is entirely parameterized by $\beta_0$.

All linear layer weights are initialized using Xavier uniform initialization 
\begin{equation}
    W \sim \mathcal{U}\left(-\sqrt{\frac{6}{n_{\text{in}} + n_{\text{out}}}}, \sqrt{\frac{6}{n_{\text{in}} + n_{\text{out}}}}\right)
\end{equation}

with biases initialized to zero. During training, the aggregated risk score $\eta(\mathbf{x})$ is optimized directly against the sampling-weighted Fine Gray partial likelihood loss. For post hoc model interpretation, each univariate shape function $f_i(x_i)$ can be visualized over its covariate support domain to inspect linear or non-linear effects on sub-distribution incidence.


\section{Loss Function}

To optimize the neural log-relative subdistribution hazard predictor under varying computational constraints, the network is trained by minimizing a sampling-aware approximation of the Fine--Gray partial likelihood. This objective integrates the additive neural predictor, accounts for right censoring through Fine--Gray risk-set weighting, and adjusts for the unequal inclusion probabilities introduced by the different risk-set sampling strategies.

The formulation begins with the Fine--Gray objective applied to the complete modified risk sets within each training mini-batch. Let $\mathcal{B}$ denote the current mini-batch, and let $\mathcal{R}_{\mathcal{B}}(T_i)$ represent the modified Fine--Gray risk set evaluated at event time $T_i$ for the target cause $k$. This risk set includes individuals who remain under observation at $T_i$, as well as those who have previously experienced a competing event.

To correctly account for the contributions of individuals retained following competing events, Inverse Probability of Censoring Weighting (IPCW) is incorporated into the Fine--Gray formulation \citep{FineGray1999}. Let $\widehat{G}(t)$ denote the Kaplan--Meier estimate of the censoring survival function, fitted using the training fold. Assuming independent right censoring, the Fine--Gray risk-set weight for an individual $j$ at event time $T_i$ is defined as:

\begin{equation}
w_{ij}^{FG}
=
\begin{cases}
1,
& T_j \geq T_i, \\[6pt]
\dfrac{\widehat{G}(T_i^-)}
{\widehat{G}(T_j^-)},
& T_j < T_i,\quad
\Delta_j \notin \{0,k\}, \\[10pt]
0,
& \text{otherwise}.
\end{cases}
\end{equation}

Under this formulation, individuals who remain under observation at the event time contribute with unit weight, while those who previously experienced a competing event receive a time-dependent adjustment. This adjustment accounts for the probability of remaining uncensored between the competing event and the current event time. Individuals who were previously censored or experienced the target event are excluded from the corresponding risk set. The formulation assumes positive censoring survival probabilities at the relevant times.

Replacing the traditional linear predictor with the additive neural risk score $\eta_{\theta,k}(\boldsymbol{X})$, the complete mini-batch Fine--Gray risk-set denominator is expressed as:

\begin{equation}
D_{\mathcal{B}}(T_i)
=
\sum_{j \in \mathcal{R}_{\mathcal{B}}(T_i)}
w_{ij}^{FG}
\exp\left(\eta_{\theta,k}(\boldsymbol{X}_j)\right),
\end{equation}

where $\theta$ represents the neural network parameters.

Let $I_k(\mathcal{B}) = \{i \in \mathcal{B} : \Delta_i = k\}$ denote the set of observed target-cause events within the mini-batch. The corresponding negative log-partial-likelihood objective is formulated as:

\begin{equation}
\mathcal{L}_{FG}(\theta)
=
-\frac{1}{|I_k(\mathcal{B})|}
\sum_{i \in I_k(\mathcal{B})}
\left[
\eta_{\theta,k}(\boldsymbol{X}_i)
-
\log D_{\mathcal{B}}(T_i)
\right].
\end{equation}

This formulation provides the unsampled reference objective against which the different risk-set sampling strategies are evaluated. Since the risk sets are constructed within individual training mini-batches rather than across the entire training cohort, the objective represents a computational approximation of the full-cohort Fine--Gray partial likelihood.

When subsampling strategies such as nested case-control, case-cohort, or the proposed dynamic entropy-guided method are employed, the complete mini-batch risk set is replaced by a sampled subset, denoted as $\widetilde{\mathcal{R}}_{\mathcal{B}}(T_i)$. However, because individuals may have unequal probabilities of inclusion, the sampled risk-set denominator must account for the resulting selection mechanism.

To address this, Inverse Probability Weighting (IPW) is applied using the inclusion probabilities associated with each sampling design \citep{HorvitzThompson1952, Prentice1986}. Let $\pi_{ij}$ denote the probability that individual $j$ is selected into the risk set associated with event $i$. The inverse sampling weight is then defined as:

\begin{equation}
w_{ij}^{S}
=
\frac{1}{\pi_{ij}}.
\end{equation}

The observed event case is always retained, such that $\pi_{ii}=1$. Individuals with lower inclusion probabilities receive larger sampling weights, compensating for their reduced representation within the sampled risk set.

The combined contribution weight for individual $j$ at event time $T_i$ is therefore obtained by multiplying the Fine--Gray censoring adjustment by the inverse sampling probability:

\begin{equation}
w_{ij}
=
w_{ij}^{FG} \cdot w_{ij}^{S}
=
\frac{w_{ij}^{FG}}{\pi_{ij}}.
\end{equation}

These two weighting components address distinct sources of adjustment. The Fine--Gray weight accounts for right censoring within the modified subdistribution risk set, whereas the sampling weight adjusts for unequal control selection probabilities.

Using these combined weights, the sampled risk-set denominator is estimated as:

\begin{equation}
\widehat{D}_{\mathcal{B}}(T_i)
=
\sum_{j \in \widetilde{\mathcal{R}}_{\mathcal{B}}(T_i)}
\frac{w_{ij}^{FG}}{\pi_{ij}}
\exp\left(\eta_{\theta,k}(\boldsymbol{X}_j)\right).
\end{equation}

The final sampling-aware negative log-partial-likelihood objective used to train the neural network is consequently defined as:

\begin{equation}
\boxed{
\widehat{\mathcal{L}}_{FG}^{IPW}(\theta)
=
-\frac{1}{|I_k(\mathcal{B})|}
\sum_{i \in I_k(\mathcal{B})}
\left[
\eta_{\theta,k}(\boldsymbol{X}_i)
-
\log\widehat{D}_{\mathcal{B}}(T_i)
\right]
}
\end{equation}

This objective allows the same neural architecture and loss formulation to be used across the different sampling strategies. For the unsampled mini-batch baseline, all eligible individuals are retained with inclusion probability equal to one, reducing the sampled objective to the complete mini-batch Fine--Gray formulation.

For numerical stability, the weighted risk-set denominator is evaluated using log-sum-exp, and small positive lower bounds are applied before taking logarithms of censoring survival estimates and sampling probabilities. Explicit upper truncation of the combined inverse-probability weights is not applied. During optimization, gradient clipping is used separately to limit extreme neural network gradients.

For fixed risk-set contributions and correctly specified positive inclusion probabilities, inverse-probability weighting recovers the complete mini-batch denominator in expectation over the sampling design. For SCC, the subcohort is drawn once and reused during training, so this design-based identity does not establish conditional unbiasedness at every optimisation step after the model has adapted to that subcohort. However, because the denominator enters the partial likelihood through a logarithmic transformation, this property does not guarantee an unbiased estimate of the complete partial-likelihood objective.

Consequently, the proposed sampling-aware loss is treated as a stochastic approximation of the complete mini-batch Fine--Gray objective. Its effectiveness is assessed empirically through predictive discrimination, calibration, training stability, and computational resource comparisons across the different sampling strategies.


% \section{Loss Function} 

% To optimize the neural relative risk predictor under varying computational constraints, the network is trained by minimizing a sampling-aware extension of the Fine-Gray partial likelihood. This objective function is mathematically designed to integrate the additive neural predictor, correct for right-censored observations, and neutralize the selection bias introduced by risk-set sampling architectures. 

% The formulation is constructed by first establishing the baseline Fine-Gray objective applied to the full cohort. Let $\mathcal{R}(T_i)$ denote the complete modified risk set evaluated at event time $T_i$ for the target cause $k$. To account for patients who are independently right-censored, Inverse Probability of Censoring Weighting (IPCW) is applied \citep{Graf1999}. If $\hat{G}(t \mid \mathbf{x})$ represents the Kaplan-Meier estimate of the censoring survival distribution, the censoring weight for a subject $j$ at time $T_i$ is computed as:
% $$ w_j^{\text{cens}}(T_i) = \frac{1}{\hat{G}(T_i^- \mid \mathbf{x}_j)} \,.$$
% Replacing the traditional linear predictor with the additive neural risk score $\eta_k(\mathbf{x})$, the baseline negative log-partial likelihood evaluated over the full risk set is formulated as:
% $$ \mathcal{L}_{\text{FG}}(\beta) = -\sum_{i:\Delta_i = k} w_i^{\text{cens}}(T_i) \left[ \eta_k(\mathbf{x}_i) - \log \sum_{j \in \mathcal{R}(T_i)} w_j^{\text{cens}}(T_i) \exp(\eta_k(\mathbf{x}_j)) \right] \,.$$

% While statistically sound, optimizing this full baseline objective across all event times is computationally exhaustive. When sub-sampling architectures (such as nested case-control or static case-cohort designs) are employed, the full risk set $\mathcal{R}(T_i)$ is replaced by a much smaller sampled subset, denoted as $\tilde{\mathcal{R}}(T_i)$. However, this subsampling distorts the natural risk distribution. To recover an unbiased representation of the original cohort, Inverse Probability Weighting (IPW) is applied based on the selection probabilities \citep{HorvitzThompson1952, Prentice1986}.

% For a given subject $j$ with an inclusion probability of $\pi_j$, the naive IPW weight is $1/\pi_j$. Because extreme fractional probabilities can destabilize neural network gradients (Introducing varince in the weights of the model), a variance reduction technique is applied by scaling the weights by the empirical marginal probability of selection, $P(S=1)$ \citep{Robins2000, Hernan2000}. This yields the stabilized sampling weight:
% $$ w_j^{\text{stab}} = \frac{P(S=1)}{\pi_j} \,.$$

% The global contribution weight for any subject $j$ evaluated at time $T_i$ is therefore the product of their stabilized sampling weight and their censoring weight:
% $$ w_{ij} = w_j^{\text{stab}} \cdot w_j^{\text{cens}}(T_i) \,.$$

% By injecting these combined weights and evaluating only over the sampled subset $\tilde{\mathcal{R}}(T_i)$, the final sampling-aware Fine-Gray objective utilized to train the model is derived as:
% $$ \mathcal{L}_{\text{FG}}^{\text{IPW}}(\beta) = -\sum_{i:\Delta_i = k} w_{ii} \left[ \eta_k(\mathbf{x}_i) - \log \sum_{j \in \tilde{\mathcal{R}}(T_i)} w_{ij} \exp(\eta_k(\mathbf{x}_j)) \right] \,,$$
% where $w_{ii}$ serves as the combined weight of the observed case at its own event time. 

% Notably, explicit weight clipping a technique frequently utilized in reinforcement learning and causal inference to bound extreme gradients is intentionally omitted from this framework. Because the sub sampled risk sets are already heavily reduced, artificially truncating valid IPW weights at arbitrary thresholds would discard critical representation information, re-introducing the very statistical bias the IPW formulation is designed to eliminate. 

% Throughout the experimental phase of this thesis, the models were aggressively optimized against this sampling-aware $\mathcal{L}_{\text{FG}}^{\text{IPW}}$ loss function. To ensure robust validation and provide a rigorous experimental control, baseline models were trained alongside them using the exact same weighted loss formulation, but evaluated comprehensively over the full, unsampled risk set.


\section{Training Details}

To evaluate performance across the datasets, the model's predictive capabilities and sampling mechanics were validated using a five-fold cross-validation framework. For the synthetic and Framingham datasets, a Stratified K-Fold splitting strategy was implemented using the event indicators. For PBC, Stratified Group K-Fold splitting used patient identifiers to keep all records belonging to a patient within the same fold. Event-aware stratification aims to preserve the proportions of target events, competing events and censored observations across folds, although exact proportions are not guaranteed, particularly when patient grouping must also be respected.

Within this cross-validation structure, all model configurations were trained using mini-batch optimisation. For each target-event time within a batch, the Full Risk (FULL) baseline evaluates all eligible observations in that batch's modified Fine--Gray risk set. The sampling strategies instead select eligible controls while retaining the current event case. This provides a comparison between full and sampled risk-set evaluation under the same mini-batch training framework.

Mini-batch training processes subsets of the training cohort, limiting the number of observations handled in each optimisation step. Within the PyTorch framework, DataLoaders supply these batches, with shuffling enabled for training and disabled for validation. Event times and event indicators are retained for each observation and used to construct the eligible Fine--Gray risk sets within each batch. When SCC, NCC or entropy-guided sampling is applied, the loss is evaluated using a further subset of eligible risk-set contributions within these batches.

Consequently, a core objective of this training protocol is to investigate whether reducing the number of risk-set contributions within each batch affects training stability, and whether the sampling-aware Fine--Gray loss with Inverse Probability Weighting (IPW) maintains convergence and predictive accuracy. This comparative analysis assesses how much additional training-time reduction risk-set sampling provides relative to the FULL baseline, and the associated trade-offs in predictive performance, with all configurations using mini-batch training.


All model configurations were optimized utilizing the Adam optimizer \citep{kingma2014adam}. Adam is highly effective in deep survival contexts as it leverages adaptive learning rates and moving averages of historical gradients to navigate the complex loss topographies generated by IPW scaling. To prevent the FeatureNets from memorizing localized noise a heightened risk when training on deeply sub-sampled risk sets L2 parameter regularization (weight decay) was integrated into the optimization loop. This regularizes the network by penalizing excessively large weights, promoting smoother feature effect curves and ensuring that the interpretable sub-networks generalize effectively to unseen clinical populations.


\section{Model Evaluation and Diagnostics}

This section details the evaluation metrics and diagnostic frameworks employed to validate the predictive performance, calibration, and computational efficiency of the competing risks models. To provide a holistic assessment, the models are evaluated on their discriminative capacity using a time-dependent Concordance Index, their absolute probabilistic accuracy via the Brier and their clinical reliability through Stratified Aalen-Johansen Cumulative Incidence Function (CIF) calibration. Furthermore, computational profiling and synthetic correlation metrics are introduced to explicitly validate the theoretical efficiency and exact feature recovery of the proposed CRISPNAM architecture.

In alignment with standard competing risk methodology, all subsequent model diagnostics, calibration curves, and feature recovery evaluations are reported exclusively for the primary event of interest (Cause 1). Competing events (Cause 2) are modeled inherently by the subdistribution architecture as competing risks, but are not the primary target of predictive evaluation.

\subsection{Brier Score}
In standard regression tasks, the Mean Squared Error (MSE) evaluates the distance between a prediction and the continuous truth. To quantify absolute probabilistic accuracy under competing risks, the Brier Score serves as the equivalent metric, formulated directly against the Cumulative Incidence Function (CIF). Evaluated at a specific clinical time horizon $t$, the Brier Score calculates the mean squared difference between the model's predicted probability of the primary event $k$ occurring prior to time $t$, and the actual observed binary outcome \citep{Graf1999}.

Because survival data is subject to right censoring meaning the exact outcome is unknown for some patients at time $t$ the naive Brier Score is heavily biased. To resolve this, an Inverse Probability of Censoring Weighting (IPCW) term, $w_i(t)$, derived from the Kaplan-Meier censoring survival function $\hat{G}(\cdot)$, is dynamically applied. For a cohort of $N$ patients, the adjusted Brier Score is defined as:
$$\text{BS}(t) = \frac{1}{N} \sum_{i=1}^N w_i(t) \left( \mathbf{1}[T_i \le t, \Delta_i = k] - \hat{F}_k(t \mid \mathbf{x}_i) \right)^2 \,,$$
where $\mathbf{1}[\cdot]$ is the indicator function denoting whether the specific event $k$ was observed prior to time $t$. A lower Brier Score reflects superior calibration and tighter probabilistic agreement with empirical reality.

% \subsection{Integrated Brier Score (IBS)}
% While the point-in-time Brier Score is useful for specific clinical horizons (e.g., the median follow-up time), it is highly dependent on the chosen evaluation point and fails to capture the model's accuracy over the entire temporal trajectory. To resolve this, the Integrated Brier Score (IBS) is utilized as a global summary metric \citep{Graf1999}. 

% The IBS computes the average of the Brier Score evaluated continuously over a prespecified temporal interval $[0, \tau]$:
% $$\text{IBS} = \frac{1}{\tau} \int_0^{\tau} \text{BS}(t) \, dt$$
% By integrating the predictive error across all valid time steps, the IBS collapses short-term and long-term probabilistic accuracy into a single scalar value. This metric is critical for objectively comparing disparate sampling strategies across the entire follow-up period.

\subsection{Concordance Index (C-Index)}
While the Brier Score measures absolute probabilistic error, the Concordance Index (C-Index) quantifies the model's discriminative ability by evaluating how accurately the predicted relative risk scores rank the observed event times. In a clinical setting, discrimination answers the question: \textit{If Patient A experiences a fatal cardiac event before Patient B, did the model assign Patient A a higher risk score?}

Let $\eta_i$ and $\eta_j$ denote the predicted log-risk scores for a pair of subjects. A pair is considered concordant if the subject who fails earlier was assigned a higher risk score ($\eta_i > \eta_j$). Because right-censoring obscures the true temporal ordering for certain patients, an IPCW adjustment is integrated to ensure unbiased evaluation. To properly evaluate non-proportional hazards where risk predictions vary over time, we rely on the time-dependent concordance index proposed by Antolini et al. \citep{Antolini2005}. Adapted for time-dependent and competing risks scenarios, the estimator is defined as:
$$\widehat{C}_{\text{IPCW}} = \frac{ \sum_{i \neq j} \mathbf{1}(T_i < T_j, \Delta_i = k) \mathbf{1}(\eta_i > \eta_j) / \hat{G}^2(T_i) }{ \sum_{i \neq j} \mathbf{1}(T_i < T_j, \Delta_i = k) / \hat{G}^2(T_i) }$$
A value of $0.5$ indicates performance equivalent to random chance, whereas a value approaching $1.0$ proves strong clinical discriminative rank-ordering.

\subsection{Stratified Aalen-Johansen CIF Calibration}
Strong discrimination (a high C-Index) guarantees that the model correctly ranks patient risk, but it does not guarantee calibration. Calibration asks: \textit{If the model predicts a 20\% risk of an event for a group of patients, do approximately 20\% of those patients actually experience the event?} To prove that the absolute risk estimations output by the Fine-Gray loss are reliable, a Stratified CIF calibration diagnostic was developed.

Under this diagnostic, the validation cohort is stratified into discrete quantiles (e.g., Low, Medium, and High risk) based on the model's predicted relative risk scores. For each stratum, the mean predicted CIF curve is plotted against empirical reality. Because the standard Kaplan-Meier estimator is mathematically invalid in the presence of competing risks, the empirical reality is estimated strictly using the non-parametric Aalen-Johansen estimator \citep{aalen1978empirical}. If the model's predicted mean CIF curve closely shadows the step-wise Aalen-Johansen empirical curve across all risk strata, it proves that the neural network's absolute risk estimates are highly calibrated.

\subsection{CPU and Memory Profiling}
A fundamental hypothesis of this thesis is that sub-sampling the risk set via static, nested, or dynamic entropy-guided designs dramatically reduces the computational burden inherent to the full cohort Fine-Gray objective. Beyond predictive accuracy, this study explicitly evaluates the computational efficiency of these sampling regimes.

A continuous resource profiler is executed alongside the cross-validation training loop. Operating on an isolated parallel thread, the profiler polls the system to record absolute CPU utilization and Resident Set Size (RSS) memory consumption, ensuring that only the resources directly consumed by the model training process are captured. By measuring resource trajectories and total execution durations independently of external background noise, the study accurately quantifies the exact computational efficiency gains unlocked by the risk set approximation methods.

\subsection{Synthetic Data Verification (Pearson \& Spearman Correlation)}
A primary limitation of analyzing real wolrd medical data is the absence of an absolute "ground truth" for individual feature effects. While clinical literature can suggest trends, the exact mathematical coefficient for a specific covariate remains fundamentally unknown. Consequently, verifying whether an interpretable model has accurately isolated a specific feature's contribution is challenging.

To bypass this limitation, the synthetic cohort experiment incorporates Pearson and Spearman rank correlation coefficients \citep{rosner2015fundamentals}. Because the synthetic dataset was generated with explicitly defined baseline hazards and exact true linear predictors, the ground truth is strictly known. The Pearson coefficient measures the degree of linear agreement between the model's predicted feature effects and the true parameters, while the Spearman coefficient evaluates the monotonic rank correlation. High correlation scores explicitly prove that the CRISPNAM architecture successfully recovers the true underlying biostatistical relationships without feature entanglement, a validation approach standard in evaluating interpretable Neural Additive Models \citep{agarwal2021neural}.

\section{Code Availability and Reproducibility}

To ensure full reproducibility of the experiments, the CRISPNAM-FG architecture, and the dynamic entropy sampling mechanisms detailed in this study, the complete Python and PyTorch codebase has been made publicly available. The open-source repository contains all data preprocessing scripts, synthetic data generation pipelines, model training configurations, and evaluation metrics. The code can be accessed and reviewed via GitHub at: 

\url{https://github.com/MichaelFalk1/Fine-Gray-CRISP-NAM-Survival-Analysis-Using-Risk-Set-Sampling.git} \citep{falk2026crispnam}.

\chapter{Results and Discussion}
\section{Dataset Characteristics}

To robustly evaluate the proposed CRISPNAM architecture, including the Fine-Gray sub-distribution loss formulation and the integrated risk-set sampling strategies, this study utilized three distinct datasets: a large-scale synthetic dataset and two smaller real wolrd clinical benchmarks (PBC and Framingham).

The synthetic dataset ($N = 15,000$) was explicitly simulated to provide a rigorous mathematical environment with a high incidence of competing events. As illustrated in Figure \ref{fig:cif_synthetic}, the pronounced divergence between the standard Kaplan-Meier estimate and the true Cumulative Incidence Function (CIF) confirms that this dataset serves as an ideal stress test. It effectively demonstrates the architecture's capacity to correct the upward bias inherent in traditional single-event survival models. Furthermore, the large sample size ensures that the nested case-control and case-cohort sampling efficiencies are validated across a wide data range.

\begin{figure}[htbp]
    \centering
    \includegraphics[width=0.75\textwidth]{images/Data/synth.png} 
    \caption{Comparison of the Cumulative Incidence Function (CIF) and Kaplan-Meier (KM) estimate for the primary event in the synthetic dataset ($N = 15,000$).}
    \label{fig:cif_synthetic}
\end{figure}

Subsequent evaluations were conducted on established, real-world benchmarks to assess the model's generalizability on standard, noisy clinical data. The Primary Biliary Cholangitis (PBC) dataset, a foundational benchmark in survival analysis, comprises 1,945 longitudinal clinical records derived from repeated patient visits across 312 unique patients. To rigorously evaluate the model on this longitudinal structure without introducing subject-level data leakage, the evaluation framework explicitly grouped these records by their respective patient identifiers. Within this structured dataset, 1,073 records (55.2\%) were right-censored, 725 (37.3\%) terminated in the primary event of interest, and 147 (7.6\%) experienced a competing event. While the lower proportion of competing events results in a less extreme deviation between the Kaplan-Meier estimate and the CIF (Figure~\ref{fig:cif_pbc}) compared to the synthetic data, the dataset remains a critical baseline for benchmarking clinical predictive performance against existing literature.

\begin{figure}[H]
    \centering
    \includegraphics[width=0.75\textwidth]{images/Data/pbc3.png} 
    \caption{Cumulative Incidence Function (CIF) versus Kaplan-Meier (KM) estimate for the primary event in the real wolrd PBC dataset.}
    \label{fig:cif_pbc}
\end{figure}

Finally, the Framingham dataset was included as a secondary real wolrd evaluation point. Though significantly smaller in scale than the synthetic cohort and heavily censored (Figure \ref{fig:cif_framingham}), it contains persistent competing risk dynamics. Utilizing this combination of datasets ensures that the model's sampling mechanisms and loss functions are comprehensively validated across both mathematically extreme theoretical conditions and standard clinical realities.

\begin{figure}[H]
    \centering
    \includegraphics[width=0.75\textwidth]{images/Data/fh2.png} 
    \caption{Cumulative Incidence Function (CIF) versus Kaplan-Meier (KM) estimate for the  Framingham dataset.}
    \label{fig:cif_framingham}
\end{figure}


\section{Optimization and Convergence}

Training deep survival architectures necessitates careful monitoring to prevent overfitting, particularly when optimizing complex sub-distribution hazard formulations across varying risk-set sampling designs. To ensure robust evaluation, all models were evaluated using a standard stratified five-fold cross-validation strategy, with a specific methodological exception applied to the Primary Biliary Cholangitis (PBC) dataset. Because the PBC data consists of repeated longitudinal measurements across 312 unique patients, we employed a stratified grouped five-fold cross-validation approach (\texttt{StratifiedGroupKFold}). This grouping mechanism ensured that all clinical visits associated with an individual patient were kept strictly within the same fold, thereby eliminating the risk of subject-level data leakage across the training and validation splits.

Given the extensive hyperparameter search space across the three datasets and four sampling strategies, plotting every trajectory is impractical. Instead, the loss curves presented reflect a single representative fold (Fold 1) utilizing the optimal hyperparameter configuration for each sampling type. The full cohort loss serves as the foundational baseline against which the sampling strategies are compared.

\subsection{Hyperparameter Stability and Selection}

The convergence stability of the CRISPNAM architecture was highly dependent on the hyperparameter configurations of the risk-set samplers. The optimal parameters were selected based on the highest cross-validated Concordance Index, revealing several consistent behavioral patterns across the datasets.

For Nested Case-Control (NCC) Sampling, the optimal number of controls was found to be $m = 5$. Increasing the number of controls did not yield significant predictive improvements but linearly increased computational overhead. This aligns with standard literature on deep survival sampling, notably Kvamme et al. \citep{kvamme2019}, who demonstrated that minimal control sets are sufficient to approximate the partial likelihood gradient effectively.

For Static Case-Cohort (SCC) Sampling, a subcohort fraction of $\alpha = 0.5$ provided the most stable optimization trajectory. While lower fractions ($\alpha \in [0.2, 0.3]$) ultimately achieved similar terminal loss values, their convergence curves were highly jagged. Extreme subcohort compression ($\alpha < 0.2$) led to gradient collapse during early epochs. Thus, $\alpha = 0.5$ is recommended for maintaining gradient stability.

The Entropy Sampling makes use of a temperature scaling parameter ($\tau$) that controls the greediness of the sampler. The optimal balance was found at $\tau = 0.5$. When the temperature was pushed toward maximum greediness ($\tau = 0.1$), the model severely overfit, as it exclusively sampled the highest-variance patients and failed to generalize to the broader distribution. Based on these empirical trajectories, a range of $\tau \in [0.5, 1.0]$ is recommended to adequately balance the exploitation of high-information risk sets with robust model generalizability.

\newpage

\subsection{Dataset Convergence Analysis}

As illustrated in Figure \ref{fig:loss_synthetic}, the Fine-Gray partial likelihood loss achieves textbook convergence on the large-scale synthetic dataset ($N=15,000$). The training and validation curves track almost identically across all four sampling configurations, verifying the structural and mathematical integrity of the gradients within the simulated competing risks environment. 

\begin{figure}[H]
    \centering
    \begin{subfigure}[b]{0.48\textwidth}
        \includegraphics[width=\textwidth]{images/Synth_results/FULL/LOSS/fold_1_loss_cause_1.png}
        \caption{Full Cohort}
    \end{subfigure}
    \hfill
    \begin{subfigure}[b]{0.48\textwidth}
        \includegraphics[width=\textwidth]{images/Synth_results/NCC/LOSS/fold_1_loss_cause_1_NCC_ncontrols_5.png}
        \caption{NCC ($m=5$)}
    \end{subfigure}
    \vskip\baselineskip
    \begin{subfigure}[b]{0.48\textwidth}
        \includegraphics[width=\textwidth]{images/Synth_results/SCC/LOSS/fold_1_loss_cause_1_SCC_subcohort_0.5.png}
        \caption{SCC ($\alpha=0.5$)}
    \end{subfigure}
    \hfill
    \begin{subfigure}[b]{0.48\textwidth}
        \includegraphics[width=\textwidth]{images/Synth_results/ENTROPY/LOSS/fold_1_loss_cause_1_Entropy_tau_0.5_compression_0.1.png}
        \caption{Entropy ($\tau=0.5$)}
    \end{subfigure}
    \caption{Loss trajectories for the synthetic dataset ($N=15,000$) under optimal hyperparameters.}
    \label{fig:loss_synthetic}
\end{figure}

The optimization trajectories for the real world PBC dataset (Figure \ref{fig:loss_pbc}) exhibit higher initial loss magnitudes, which is characteristic of standard, unscaled clinical data. However, the model successfully minimizes the sub-distribution hazard across all sampling configurations without catastrophic divergence. The validation loss consistently mirrors the training loss, proving that the network's internal regularization mechanisms prevent it from memorizing the sampled clinical risk-sets.

\begin{figure}[H]
    \centering
    \begin{subfigure}[b]{0.48\textwidth}
        \includegraphics[width=\textwidth]{images/pbc_results/FULL/LOSS/fold_1_loss_cause_1.png}
        \caption{Full Cohort}
    \end{subfigure}
    \hfill
    \begin{subfigure}[b]{0.48\textwidth}
        \includegraphics[width=\textwidth]{images/pbc_results/NCC/LOSS/fold_1_loss_cause_1_NCC_ncontrols_5.png}
        \caption{NCC ($m=5$)}
    \end{subfigure}
    \vskip\baselineskip
    \begin{subfigure}[b]{0.48\textwidth}
        \includegraphics[width=\textwidth]{images/pbc_results/SCC/LOSS/fold_1_loss_cause_1_SCC_subcohort_0.5.png}
        \caption{SCC ($\alpha=0.5$)}
    \end{subfigure}
    \hfill
    \begin{subfigure}[b]{0.48\textwidth}
        \includegraphics[width=\textwidth]{images/pbc_results/ENTROPY/LOSS/fold_1_loss_cause_1_Entropy_tau_0.5_compression_0.1.png}
        \caption{Entropy ($\tau=0.5$)}
    \end{subfigure}
    \caption{PBC dataset loss trajectories demonstrating stable convergence across all sampling methods.}
    \label{fig:loss_pbc}
\end{figure}

Finally, the convergence trajectories for the  Framingham dataset (Figure \ref{fig:loss_fh}) present a unique visual scale. Because the initial unoptimized training loss is exceptionally high across all sampling configurations, the y-axis is heavily compressed. This visual compression makes the validation loss appear entirely flat at first glance. However, numerical inspection of the optimization logs confirms that the validation loss does indeed decrease steadily (e.g., dropping from approximately 6.6 to 6.1) as the network optimizes. 

This confirms that the model is actively learning and generalizing the sub-distribution hazard rather than stagnating. Importantly, across Full Cohort, NCC, SCC, and Entropy sampling, the training loss converges smoothly down to effectively meet the validation loss without crossing into severe overfitting territory. This further validates the robustness of the chosen hyperparameters when navigating heavily censored competing risks data.

\begin{figure}[H]
    \centering
    \begin{subfigure}[b]{0.48\textwidth}
        \includegraphics[width=\textwidth]{images/fh_results/FULL/LOSS/fold_1_loss_cause_1.png}
        \caption{Full Cohort}
    \end{subfigure}
    \hfill
    \begin{subfigure}[b]{0.48\textwidth}
        \includegraphics[width=\textwidth]{images/fh_results/NCC/LOSS/fold_1_loss_cause_1_NCC_ncontrols_5.png}
        \caption{NCC ($m=5$)}
    \end{subfigure}
    \vskip\baselineskip
    \begin{subfigure}[b]{0.48\textwidth}
        \includegraphics[width=\textwidth]{images/fh_results/SCC/LOSS/fold_1_loss_cause_1_SCC_subcohort_0.5.png}
        \caption{SCC ($\alpha=0.5$)}
    \end{subfigure}
    \hfill
    \begin{subfigure}[b]{0.48\textwidth}
        \includegraphics[width=\textwidth]{images/fh_results/ENTROPY/LOSS/fold_1_loss_cause_1_Entropy_tau_0.5_compression_0.1.png}
        \caption{Entropy ($\tau=0.5$)}
    \end{subfigure} 
    \caption{Framingham dataset loss trajectories showing stable convergence despite visual y-axis compression.}
    \label{fig:loss_fh}
\end{figure}

\subsection{Computational Resource Profiling}

A primary motivation for implementing risk-set sampling strategies in deep survival architectures is the reduction of computational overhead. Calculating the exact sub-distribution partial likelihood for the full cohort becomes computationally intractable as dataset sizes scale. To quantify this efficiency gain, hardware resource utilization (CPU and Memory) and execution times were profiled during model training. 

Given that the computational patterns were consistent across all experiments, the resource profiles are visualized exclusively for the large-scale synthetic dataset ($N=15,000$), which serves as the upper-bound stress test for the architecture. The smaller PBC and Framingham datasets exhibited identical behavioral patterns at lower absolute magnitudes.

\subsubsection{Resource Utilization and Execution Efficiency}

Across all sampling configurations, CPU utilization remained highly stable at approximately 50\% during the active training phase, indicating consistent dataloader throughput without severe processing bottlenecks. 

Memory utilization exhibited a characteristic pattern across all methods: an initial allocation spike during data ingestion and model initialization, followed by a stable plateau during batched gradient updates. 

The baseline full cohort calculation required approximately 725 MB during initialization and plateaued near 600 MB. However, it required the longest execution time, taking roughly 60 seconds per fold iteration.

Both SCC ($\alpha=0.5$) and NCC ($m=5$) required slightly higher initial memory spikes (approx. 790 MB and 780 MB, respectively) due to the overhead of the sampling algorithms allocating the risk-sets in memory. However, by compressing the active risk pool, both methods reduced the execution time to approximately 45 seconds a 25\% reduction in compute time compared to the full cohort.

The entropy sampler ($\tau=0.5$) exhibited the highest initialization memory peak (approx. 830 MB), reflecting the upfront computational cost of calculating the entropy distribution across the risk-sets. Strikingly, once initialized, this strategy proved to be the most efficient, completing the execution in roughly 28 seconds. This represents a greater than 50\% reduction in training time compared to the baseline, proving that intelligently exploiting high-information risk sets drastically accelerates optimization.

Ultimately, these profiles confirm that while advanced sampling strategies require marginally more initialization memory to construct the active cohorts, they heavily optimize the overall computational pipeline by drastically reducing the time required for gradient convergence.

\section{Predictive Performance and Training Time}
\subsection{Synthetic Dataset}

The synthetic dataset ($N=15,000$) provides a robust mathematical stress test featuring a high incidence of competing events. This environment enables a clear evaluation of the CRISPNAM architecture's capacity to resolve competing risk bias while establishing the baseline predictive performance and computational overhead required for full partial likelihood calculations.

\begin{figure}[H]
    \centering
    \includegraphics[width=0.85\textwidth]{images/Synth_results/Pooled_C-Index.png}
    \caption{Pooled C-Index scores across all sampling strategies on the synthetic dataset.}
    \label{fig:synth_c_index}
\end{figure}


The full cohort model serves as the theoretical and empirical baseline. Because it utilizes the entire risk-set at every event time, it guarantees the most precise gradient updates at the expense of maximum computational complexity. As illustrated in Figures \ref{fig:synth_c_index} and \ref{fig:synth_brier_score}, the full cohort achieved a pooled C-index of $0.850$ and a Brier Score of $0.107$. However, this mathematical precision comes with a significant computational penalty. The full cohort training required an average of $58.16 \pm 21.87$ seconds per fold, establishing the baseline latency that the risk-set sampling strategies aim to optimize.

\begin{figure}[H]
    \centering
    \includegraphics[width=0.85\textwidth]{images/Synth_results/Training_Duration_mean_std_error.png}
    \caption{Average training duration per fold on the synthetic dataset.}
    \label{fig:synth_training_duration}
\end{figure}



The performance and computational cost of the NCC sampler are highly dependent on the number of controls ($m$). When restricted to $m=5$, the model experiences a slight degradation in discriminative capability (C-index: $0.835$) and calibration (Brier Score: $0.139$). Expanding the control pool to $m=100$ restores performance to baseline parity (C-index: $0.849$, Brier Score: $0.107$). Crucially, however, as the number of controls increases, both the training duration and the cross-fold variance scale linearly, with $m=100$ actually exceeding the training time of the full cohort (approaching 85 seconds on average). This behavior mirrors findings documented by \cite{kvamme2019} for deep single-event Cox models, where arbitrarily large control sets introduce severe computational bottlenecks without proportional predictive gains. Furthermore, the massive standard deviation bars at higher control limits indicate that standard NCC struggles with execution stability across randomized data splits.

\begin{figure}[H]
    \centering
    \includegraphics[width=0.85\textwidth]{images/Synth_results/Pooled_Brier.png}
    \caption{Pooled Brier Scores on the synthetic dataset.}
    \label{fig:synth_brier_score}
\end{figure}


The Static Case-Cohort sampler demonstrated high stability across varying subcohort fractions ($\alpha$). Even under extreme compression ($\alpha=0.05$), the model maintained a C-index of $0.842$ and a Brier Score of $0.110$. At the optimal configuration of $\alpha=0.5$, the metrics aligned near-perfectly with the full cohort (C-index: $0.849$, Brier Score: $0.107$) while compressing the training time to roughly 38 seconds. The tighter variance bars observed at $\alpha=0.5$ indicate that maintaining a static background cohort effectively stabilizes the model's convergence timing compared to the highly variable NCC approach.


The Entropy sampler provided the most significant finding of the synthetic evaluation. At optimal temperature scaling ($\tau \in [0.5, 1.0]$), the entropy sampler not only matched the discriminative performance of the Full Cohort baseline (C-index: $0.851$ at $\tau=0.5$) and maintained parity in calibration (Brier Score: $0.107$), but it achieved this in a fraction of the time. 

The training duration for the entropy sampler at $\tau=0.5$ was drastically reduced to approximately 29 seconds a 50\% reduction compared to the full cohort baseline. Furthermore, the cross-fold variance for training duration in the $\tau \in [0.5, 1.0]$ range is exceptionally small. This negligible standard deviation proves that the entropy formulation is highly deterministic; by intelligently weighting high-information observations and filtering out redundant background noise, the sampler accelerates convergence without sacrificing predictive stability.

\subsection{Real World Clinical Benchmark Primary Biliary Cholangitis (PBC) Dataset}

While the synthetic dataset validated the mathematical limits of the architecture under severe competing risk conditions, the Primary Biliary Cirrhosis (PBC) dataset serves to evaluate the model's robustness against standard clinical noise, unscaled covariates, and a smaller overall sample size. 

Because the dataset is significantly smaller, the baseline execution time is inherently lower. However, the performance metrics presented in Figures \ref{fig:pbc_c_index}, \ref{fig:pbc_brier_score}, and \ref{fig:pbc_training_duration} reveal that the sampling strategies react far more violently to hyperparameter tuning on noisy clinical data than they do on clean synthetic data.

\begin{figure}[H]
    \centering
    \includegraphics[width=0.85\textwidth]{images/pbc_results/Pooled_C-Index_PBC.png}
    \caption{Pooled C-Index scores across sampling strategies on the clinical PBC dataset.}
    \label{fig:pbc_c_index}
\end{figure}

The full cohort model maintained strong discriminative capabilities on the clinical data, achieving a pooled C-index of 0.756 and a Brier Score of 0.180. Due to the smaller dataset size, the computational burden was relatively light, averaging roughly 28 seconds per fold. Consequently, the primary objective of risk-set sampling on the PBC dataset shifts slightly: rather than strictly minimizing execution time, the goal is to prove that sampling does not catastrophically degrade the predictive integrity of a smaller cohort.

\begin{figure}[H]
    \centering
    \includegraphics[width=0.85\textwidth]{images/pbc_results/Pooled_Brier_PBC.png}
    \caption{Pooled Brier Scores on the PBC dataset.}
    \label{fig:pbc_brier_score}
\end{figure}

The NCC sampler exhibited somewhat volatile but competitive behavior on the clinical data. As the control parameter increased from $m=5$ to $m=10$, the predictive performance jumped from a C-index of 0.666 to a peak of 0.779, outperforming the full cohort, before stabilizing around 0.772 to 0.774 for larger control sets up to $m=100$. However, the execution time scaled linearly alongside the control size. At $m=100$, NCC required approximately 34 seconds, exceeding the duration of the full cohort model. This reinforces the finding that arbitrarily expanding control sets is computationally inefficient; an optimal balance for NCC on this dataset is found at lower control sizes like $m=10$, which maximizes discriminative performance (Brier score of 0.168) while minimizing training time.

\begin{figure}[H]
    \centering
    \includegraphics[width=0.85\textwidth]{images/pbc_results/Training_Duration_mean_std_error_PBC.png}
    \caption{Average training duration per fold on the PBC dataset.}
    \label{fig:pbc_training_duration}
\end{figure}

The limitations of the Static Case-Cohort (SCC) sampler were starkly exposed on the clinical data. At extreme subcohort compressions ($\alpha=0.05$), the model's discriminative ability collapsed entirely, plunging to a C-index of 0.338, performing worse than random chance. The calibration similarly degraded, yielding a Brier Score of 0.285. This catastrophic failure indicates that when the overall dataset is small and noisy, heavily compressing the static background cohort removes critical variance necessary for stable gradient descent. The SCC model only approaches viability at larger subcohort fractions like $\alpha=0.3$ and $\alpha=0.5$ (achieving C-indices of 0.730 and 0.733, respectively), but it consistently fails to match the baseline full cohort performance.

The introduction of clinical noise forced a critical shift in the optimal temperature parameter ($\tau$) for the Entropy sampler. While greedy sampling might excel on clean data, applying a highly greedy temperature ($\tau=0.5$ or $\tau=1.0$) to the PBC data resulted in suboptimal discriminative performance (C-indices of 0.734 and 0.727) and poor calibration, such as a Brier Score of 0.209 for $\tau=0.5$. 

This occurs because a low $\tau$ forces the sampler to exclusively focus on the highest-variance observations. In a noisy clinical setting, these "high-variance" patients are often extreme outliers. By oversampling outliers, the model fails to learn the general underlying survival distribution. 

The architecture recovers brilliantly when the temperature is relaxed to $\tau=5.0$ or $\tau=10.0$. By decreasing the greediness and introducing a flatter probability distribution over the risk-sets, the Entropy sampler recovers C-indices of 0.753 and 0.752 respectively, reaching statistical parity with the full cohort's 0.756. The calibration also strictly improves in this relaxed regime, achieving a Brier Score of 0.174 at $\tau=10.0$. Furthermore, at the optimal balance of $\tau=5.0$, the execution time remains highly efficient at approximately 16 seconds (nearly half the time of the full cohort model), confirming that Entropy sampling remains a superior optimization strategy provided the temperature is scaled appropriately to the dataset's intrinsic noise profile.

\subsection{Real World Framingham Dataset}

The Framingham dataset introduces a final, unique environmental challenge: it is heavily right-censored and features a lower overall event density. This scenario tests how the CRISPNAM architecture and its sampling mechanisms handle high uncertainty within the active risk pool. The performance and computational profiling for this dataset are presented in Figures \ref{fig:fh_c_index}, \ref{fig:fh_brier_score}, and \ref{fig:fh_training_duration}.

\begin{figure}[H]
    \centering
    \includegraphics[width=0.85\textwidth]{images/fh_results/Pooled_C-Index_FH.png}
    \caption{Pooled C-Index scores on the highly censored Framingham dataset.}
    \label{fig:fh_c_index}
\end{figure}


Due to the heavy censoring and inherent noise, the Full Cohort baseline achieves a modest pooled C-index of 0.645 and a high  Brier Score of 0.670. Because the dataset is small, the entire risk matrix easily fits into GPU memory, allowing highly optimized, vectorized tensor operations. Consequently, the Full Cohort training is executed rapidly, averaging roughly 21 seconds per fold. 

\begin{figure}[H]
    \centering
    \includegraphics[width=0.85\textwidth]{images/fh_results/Pooled_Brier_FH.png}
    \caption{Pooled Brier Scores on the Framingham dataset.}
    \label{fig:fh_brier_score}
\end{figure}


The Framingham evaluation reveals a counter-intuitive but critical behavior of the NCC sampler: it improves predictive performance but severely degrades computational efficiency. Across all control sizes ($m \in [5, 100]$), NCC consistently outperforms the Full Cohort baseline, achieving a C-index of up to 0.677 ($m=50$) and lowering the Brier Score to 0.644 ($m=5$). In heavily censored datasets, utilizing the full cohort forces the model to optimize against a massive pool of noisy, uncertain observations. NCC mitigates this by randomly discarding background noise, effectively acting as a data-level regularizer that prevents overfitting.

However, the computational cost is inverted. At $m=5$, NCC requires approximately 34 seconds to train a 60\% increase over the Full Cohort. This is believed to occur because the algorithmic overhead of dynamically sampling controls via CPU-bound loops outpaces the GPU compute time saved by the reduced matrix size. This confirms that on small datasets, the vectorization benefits of the Full Cohort outweigh the data-compression benefits of NCC.

\begin{figure}[H]
    \centering
    \includegraphics[width=0.85\textwidth]{images/fh_results/Training_Duration_mean_std_error_FH.png}
    \caption{Average training duration per fold on the Framingham dataset.}
    \label{fig:fh_training_duration}
\end{figure}


Similar to the PBC dataset, SCC struggles with stability when the background cohort is too small. At $\alpha=0.05$, the C-index collapses to 0.588. While it stabilizes at $\alpha=0.5$ (C-index: 0.637), it fails to capture the regularization benefits seen in NCC and suffers from high calibration error (Brier Score: 0.681). Its execution time remains comparable to the baseline, offering no distinct advantage in this highly censored environment.


The Entropy sampler reinforces the tuning dynamics observed in the PBC dataset. At an aggressively greedy temperature ($\tau=0.1$), the model overfits to the highest variance clinical outliers, resulting in a collapsed C-index of 0.575 and a severe Brier Score spike (0.863). 

However, when properly tuned to $\tau=5.0$, the entropy mechanism successfully filters the heavy censoring noise. It achieves a C-index of 0.663 and a Brier Score of 0.655, outperforming the Full Cohort baseline. Furthermore, unlike NCC, the algorithmic implementation of the Entropy sampler is highly efficient at optimal temperatures; at $\tau=0.5$, it executes in roughly 17 seconds, successfully improving upon both the predictive and computational benchmarks of the Full Cohort.


\section{Model Calibration and Discrimination}

While the previous section established the predictive performance using the pooled Concordance Index and Brier Score, these metrics primarily evaluate general rank-ordering and time-averaged error. To fully validate the CRISPNAM architecture, it is necessary to examine its calibration its ability to predict absolute event probabilities accurately over time and its capacity to recover true underlying biostatistical relationships.

\subsection{Ground Truth Risk Recovery}

A fundamental limitation of real wolrd clinical datasets is that the true feature effects and baseline hazards are strictly unknown. The synthetic dataset ($N=15,000$) bypasses this limitation, allowing for a direct comparison between the model's predicted relative risk score ($\hat{\eta}$) and the mathematically defined True Linear Predictor.

Figure \ref{fig:synth_truth_recovery} illustrates this relationship for the Full Cohort baseline. The architecture successfully isolates the true hazard contributions, achieving a Spearman rank correlation of 0.833 and a Pearson linear correlation of 0.843 for the primary event of interest. This strong linear and monotonic agreement proves that the deep neural additive framework accurately recovers the true biostatistical feature effects without entanglement, even under severe competing risk conditions.

\begin{figure}[htbp]
    \centering
    \includegraphics[width=0.85\textwidth]{images/Synth_results/FULL/Metrics/model_diagnostics_cause_1.png}
    \caption{Model diagnostics for the Full Cohort on the synthetic dataset, detailing rank correlation and subdistribution trajectory bounds.}
    \label{fig:synth_truth_recovery}
\end{figure}

\subsection{Cumulative Incidence Calibration}

Beyond rank-ordering patients correctly, a survival model must be highly calibrated; if the network predicts a 40\% cumulative incidence for a specific subgroup at a specific time horizon, empirical reality must reflect that exact proportion. 

To validate this absolute accuracy, the synthetic cohort was partitioned into three distinct risk strata (Low, Medium, and High) based on the network's predicted scores. As shown in the right panel of Figure \ref{fig:synth_truth_recovery}, the predicted mean CIF curves for all three strata effectively overlay the non-parametric Empirical Aalen-Johansen estimators. There is no visible divergence, probability collapse, or inflation at the temporal bounds, confirming that the unmodified Fine-Gray objective function is deeply calibrated.

\begin{figure}[htbp]
    \centering
    \includegraphics[width=0.85\textwidth]{images/Synth_results/FULL/Metrics/subject_curves_cause_1_cif.png}
    \caption{Individual patient Cumulative Incidence trajectories compared against the population Aalen-Johansen baseline (Synthetic Dataset).}
    \label{fig:synth_individual_cif}
\end{figure}

\subsection{Entropy Sampler Calibration and Risk Recovery}

A core hypothesis of this study is that intelligent risk-set sampling can drastically reduce computational overhead without corrupting the model's calibration or its ability to recover true biostatistical relationships. To evaluate this, the dynamic Entropy sampling strategy was benchmarked against the Full Cohort baseline.

Remarkably, the diagnostic evaluations reveal that the Entropy sampler does not merely match the baseline; at optimal temperature scaling ($\tau=0.5$), it strictly outperforms it in ground truth recovery. As illustrated in Figure \ref{fig:entropy_truth_recovery}, the Entropy sampler achieves a Pearson linear correlation of 0.866 and a Spearman rank correlation of 0.857. This is a measurable improvement over the Full Cohort's respective scores of 0.843 and 0.833. By dynamically prioritizing high-variance, high-information subjects and aggressively down-sampling dense, redundant patient clusters, the entropy mechanism effectively acts as a data-level regularizer, allowing the neural additive framework to isolate a cleaner underlying hazard signal.

Furthermore, as the temperature parameter $\tau$ is increased forcing the probability distribution to become flatter and less greedy the model's recovery metrics gracefully degrade back toward the Full Cohort baseline. At $\tau=1.0$, the Pearson correlation drops to 0.853, and at $\tau=5.0$, it falls to 0.839. This confirms the expected theoretical behavior: as the dynamic sampler approaches a uniform random sample, its regularizing benefits diminish.

\begin{figure}[H]
    \centering
    \includegraphics[width=0.85\textwidth]{images/Synth_results/ENTROPY/Metrics/model_diagnostics_cause_1_Entropy_tau_0.5_compression_0.1.png}
    \caption{Model diagnostics for the optimal Entropy Sampler ($\tau=0.5$) demonstrating ground-truth risk recovery and IPW-adjusted calibration.}
    \label{fig:entropy_truth_recovery}
\end{figure}

Crucially, this aggressive sub-sampling does not destroy the absolute probability estimates. A known vulnerability of risk-set sampling is the introduction of severe selection bias. However, the right panel of Figure \ref{fig:entropy_truth_recovery} demonstrates that the Inverse Probability Weighting (IPW) integration perfectly neutralizes this bias. The predicted Stratified CIF curves cleanly track the empirical Aalen-Johansen estimates without probability inflation.

\begin{figure}[H]
    \centering
    \includegraphics[width=0.85\textwidth]{images/Synth_results/ENTROPY/Metrics/subject_curves_cause_1_cif_Entropy_tau_0.5_compression_0.1.png}
    \caption{Individual patient Cumulative Incidence trajectories generated via the Entropy sampler ($\tau=0.5$).}
    \label{fig:entropy_individual_cif}
\end{figure}

This calibration is further maintained at the micro-level. Figure \ref{fig:entropy_individual_cif} displays the individual subject CIF trajectories generated by the $\tau=0.5$ Entropy model. Despite being trained on a heavily fractionated approximation of the cohort, the network continues to output stable, monotonic, and bounded incidence curves that appropriately scale to specific patient covariate profiles.


\subsection{Nested Case-Control (NCC) Calibration and Risk Recovery}

Following the evaluation of the dynamic Entropy sampler, the Nested Case-Control (NCC) architecture was assessed to determine how rigid, uniform control sampling impacts absolute calibration and feature recovery. As established in the macro-performance analysis, NCC is highly sensitive to the defined control quota ($m$).

When provided with a large control pool ($m=100$), the NCC sampler successfully mimics the full cohort. As shown in Figure \ref{fig:ncc_truth_recovery_100}, this configuration achieves a Pearson correlation of 0.856 and a Spearman rank correlation of 0.847, marginally outperforming the Full Cohort baseline by filtering out a portion of background noise. The Inverse Probability Weighting (IPW) mechanism functions effectively at this scale, ensuring the predicted stratified CIF curves accurately align with the empirical Aalen-Johansen estimates without significant bias.

\begin{figure}[htbp]
    \centering
    \includegraphics[width=0.85\textwidth]{images/Synth_results/NCC/Metrics/model_diagnostics_cause_1_NCC_ncontrols_100.png}
    \caption{Model diagnostics for NCC with an optimal control pool ($m=100$).}
    \label{fig:ncc_truth_recovery_100}
\end{figure}

However, evaluating the boundary conditions of the NCC sampler reveals a critical vulnerability when the active risk-set is overly compressed. Figure \ref{fig:ncc_truth_recovery_5} illustrates the catastrophic breakdown of the model when restricted to just $m=5$ controls per case. While the linear correlation degrades to 0.806, the most severe failure occurs in the absolute probability calibration. 

\begin{figure}[H]
    \centering
    \includegraphics[width=0.85\textwidth]{images/Synth_results/NCC/Metrics/model_diagnostics_cause_1_NCC_ncontrols_5.png}
    \caption{Model diagnostics for a heavily constrained NCC sampler ($m=5$), demonstrating structural calibration failure in the higher risk strata.}
    \label{fig:ncc_truth_recovery_5}
\end{figure}

The Stratified CIF calibration plot for $m=5$ demonstrates a total structural collapse. The high-risk curve (Stratum 3) terminates prematurely near $t=100$, and all strata drastically underestimate the true Aalen-Johansen incidence. This occurs because starving the neural network of valid controls artificially inflates the perceived baseline hazard during mini-batch optimization, distorting the IPW scaling. This visual diagnostic directly explains why the Brier Score spiked to 0.139 in the previous section: without sufficient controls, standard NCC cannot maintain absolute probabilistic accuracy over time.

\begin{figure}[htbp]
    \centering
    \includegraphics[width=0.85\textwidth]{images/Synth_results/NCC/Metrics/subject_curves_cause_1_cif_NCC_ncontrols_100.png}
    \caption{Individual patient Cumulative Incidence trajectories generated via NCC ($m=100$).}
    \label{fig:ncc_individual_cif}
\end{figure}

Conversely, when operating at the optimal $m=100$ configuration, the micro-level calibration remains intact. Figure \ref{fig:ncc_individual_cif} confirms that the individual patient trajectories are stable, monotonic, and effectively separate high-risk and low-risk subjects while respecting the mathematical boundaries of the subdistribution hazard framework.

\subsection{Static Case-Cohort (SCC) Calibration and Risk Recovery}

The Static Case-Cohort (SCC) architecture offers a fundamentally different sampling philosophy, relying on a fixed background subcohort rather than time-matched controls. While this approach provides stable optimization across iterations, its ability to maintain absolute calibration is highly dependent on the retention fraction ($\alpha$).

When provided with a robust background cohort ($\alpha=0.5$), the SCC sampler demonstrates strong structural integrity. As shown in Figure \ref{fig:scc_truth_recovery_05}, this configuration recovers the ground truth linear predictor with a Pearson correlation of 0.849 and a Spearman correlation of 0.840, slightly exceeding the Full Cohort baseline. Importantly, the IPW integration successfully scales the fixed background cohort to match the population distribution, resulting in Stratified CIF curves that accurately align with the empirical Aalen-Johansen estimates across all three risk strata.

\begin{figure}[H]
    \centering
    \includegraphics[width=0.85\textwidth]{images/Synth_results/SCC/Metrics/model_diagnostics_cause_1_SCC_subcohort_0.5.png}
    \caption{Model diagnostics for SCC with an optimal subcohort fraction ($\alpha=0.5$).}
    \label{fig:scc_truth_recovery_05}
\end{figure}

However, similar to the NCC architecture, aggressively compressing the risk-set exposes severe methodological limits. When the background subcohort is reduced to just 5\% of the total population ($\alpha=0.05$), the model's calibration begins to fracture, as illustrated in Figure \ref{fig:scc_truth_recovery_05_collapse}. 

\begin{figure}[htbp]
    \centering
    \includegraphics[width=0.85\textwidth]{images/Synth_results/SCC/Metrics/model_diagnostics_cause_1_SCC_subcohort_0.05.png}
    \caption{Model diagnostics for a heavily compressed SCC sampler ($\alpha=0.05$) highlighting the degradation of IPW scaling.}
    \label{fig:scc_truth_recovery_05_collapse}
\end{figure}

Interestingly, the nature of the SCC collapse differs from the NCC collapse. While the severely constrained NCC model ($m=5$) completely lost structural integrity and terminated early, the heavily constrained SCC model ($\alpha=0.05$) maintains its basic shape and still achieves a reasonable Pearson correlation of 0.835. The failure occurs entirely in the absolute calibration: because the static subcohort is too sparse to accurately represent the population's covariate density, the IPW scaling becomes skewed. Consequently, the predicted Stratified CIF curves systematically pull away from the true empirical baseline, particularly for the low and medium risk groups.

\begin{figure}[htbp]
    \centering
    \includegraphics[width=0.85\textwidth]{images/Synth_results/SCC/Metrics/subject_curves_cause_1_cif_SCC_subcohort_0.5.png}
    \caption{Individual patient Cumulative Incidence trajectories generated via SCC ($\alpha=0.5$).}
    \label{fig:scc_individual_cif}
\end{figure}

Provided the subcohort fraction remains sufficiently large (e.g., $\alpha=0.5$), the SCC architecture remains highly reliable at the individual patient level. Figure \ref{fig:scc_individual_cif} demonstrates that the optimal SCC configuration successfully generates individualized, monotonic CIF trajectories that respect the subdistribution boundaries and cleanly separate high and low-risk clinical profiles.


\section{Real World Model Calibration and Discrimination}

While the synthetic dataset provided a mathematically pure environment to test absolute boundaries and feature recovery, standard clinical data introduces high variance, unscaled covariates, and inherent observational noise. To validate the CRISPNAM architecture's generalizability, the calibration of the risk-set sampling mechanisms was evaluated against the real wolrd Primary Biliary Cirrhosis (PBC) benchmark.

Because the ground truth linear predictor is fundamentally unknown in real wolrd clinical data, the evaluation relies exclusively on empirical calibration diagnostics. Specifically, the model's predicted Cumulative Incidence Functions (CIF) are plotted against the non-parametric Aalen-Johansen estimators, evaluated both at a stratified group level and via marginal calibration gaps.

\subsection{Real World Clinical Benchmark Primary Biliary Cholangitis Dataset}

\subsubsection{Full Cohort Baseline}

The Full Cohort configuration serves as the empirical baseline for the PBC dataset. By computing the partial likelihood over the entire available risk set at every event time, it establishes the maximum achievable calibration given the intrinsic noise of the data.

Figure \ref{fig:pbc_full_calibration} illustrates this baseline calibration. The left panel displays the Stratified CIF Calibration; the clinical cohort was partitioned into three distinct risk strata (Low, Medium, and High) based on the network's predicted relative risk scores. The model's predicted mean incidence (solid lines) tracks the empirical Aalen-Johansen estimates (dashed lines) closely across all three risk levels. While there is minor localized deviation a standard artifact of noisy, real wolrd medical data the overall probabilistic scale is accurate, and there is no severe divergence or probability collapse over the 5,000-day observational window.

\begin{figure}[htbp]
    \centering
    \begin{subfigure}[b]{0.48\textwidth}
        \includegraphics[width=\textwidth]{images/pbc_results/FULL/Metrics/cif_calibration_cause_1.png}
        \caption{Stratified CIF Calibration}
    \end{subfigure}
    \hfill
    \begin{subfigure}[b]{0.48\textwidth}
        \includegraphics[width=\textwidth]{images/pbc_results/FULL/Metrics/subject_curves_cause_1_cif.png}
        \caption{Marginal Calibration and Discrimination}
    \end{subfigure}
    \caption{Calibration diagnostics for the Full Cohort baseline on the PBC dataset.}
    \label{fig:pbc_full_calibration}
\end{figure}

The right panel of Figure \ref{fig:pbc_full_calibration} introduces the Marginal Calibration diagnostic. This metric compares the mean of all individual patient predictions (orange line) directly against the overall population Aalen-Johansen baseline (dashed black line). The shaded "Calibration Gap" remains exceptionally narrow throughout the entire timeline, confirming that the Full Cohort model yields unbiased absolute risk estimates without systematically over- or under-predicting the incidence of the primary event.

\subsubsection{Nested Case-Control (NCC)}

The calibration of the Nested Case-Control (NCC) architecture on the PBC dataset heavily mirrors the failure states observed in the synthetic environment, proving that the algorithm's vulnerability to extreme risk-set compression is a structural limitation, not an artifact of data simulation.

When provided with a robust control pool ($m=100$), the NCC sampler successfully recovers the baseline calibration. As seen in Figure \ref{fig:pbc_ncc_100_calibration}, both the Stratified CIF and Marginal Calibration plots look nearly identical to the Full Cohort. The IPW mechanism functions as intended, neutralizing the sampling bias and preventing the Marginal Calibration Gap from expanding.

\begin{figure}[H]
    \centering
    \begin{subfigure}[b]{0.48\textwidth}
        \includegraphics[width=\textwidth]{images/pbc_results/NCC/Metrics/cif_calibration_cause_1_NCC_ncontrols_100.png}
        \caption{Stratified CIF Calibration ($m=100$)}
    \end{subfigure}
    \hfill
    \begin{subfigure}[b]{0.48\textwidth}
        \includegraphics[width=\textwidth]{images/pbc_results/NCC/Metrics/subject_curves_cause_1_cif_NCC_ncontrols_100.png}
        \caption{Marginal Calibration ($m=100$)}
    \end{subfigure}
    \caption{Calibration diagnostics for NCC with a large control pool ($m=100$) on the PBC dataset.}
    \label{fig:pbc_ncc_100_calibration}
\end{figure}

However, when the NCC sampler is aggressively constrained to just $m=5$ controls per case, the absolute calibration degrades significantly, despite the dataset being much smaller than the synthetic cohort. Figure \ref{fig:pbc_ncc_5_calibration} illustrates this breakdown. In the left panel, the high-risk stratum (green line) systematically underestimates the empirical reality, pulling away from the Aalen-Johansen estimator early in the timeline. 

\begin{figure}[htbp]
    \centering
    \begin{subfigure}[b]{0.48\textwidth}
        \includegraphics[width=\textwidth]{images/pbc_results/NCC/Metrics/cif_calibration_cause_1_NCC_ncontrols_5.png}
        \caption{Stratified CIF Calibration ($m=5$)}
    \end{subfigure}
    \hfill
    \begin{subfigure}[b]{0.48\textwidth}
        \includegraphics[width=\textwidth]{images/pbc_results/NCC/Metrics/subject_curves_cause_1_cif_NCC_ncontrols_5.png}
        \caption{Marginal Calibration ($m=5$)}
    \end{subfigure}
    \caption{Calibration diagnostics for a heavily constrained NCC sampler ($m=5$) detailing the resulting calibration gap expansion.}
    \label{fig:pbc_ncc_5_calibration}
\end{figure}

The consequence of this failure is visually stark in the Marginal Calibration plot (right panel). The mean model prediction severely detaches from the population baseline, creating a massive, expanding Calibration Gap. By starving the neural network of valid controls at each event time, the IPW formulation loses its stabilizing context, causing the model to systematically underestimate the baseline hazard of the clinical population. This definitively proves that while ultra-low control NCC ($m=5$) provides maximum computational speed, it sacrifices the absolute probabilistic integrity required for reliable clinical deployment.

\subsubsection{Static Case-Cohort (SCC)}

The Static Case-Cohort (SCC) architecture relies on a fixed background subcohort, theoretically providing a more stable optimization target than the dynamically changing risk sets of the NCC approach. However, applying this method to the small, noisy Primary Biliary Cirrhosis (PBC) dataset reveals significant vulnerabilities to cohort sparsity.

Even under an optimal, low-compression configuration ($\alpha=0.5$), where 50\% of the population is retained as the background cohort, the SCC sampler struggles to match the baseline calibration. As shown in Figure \ref{fig:pbc_scc_50_calibration}, while the general shape of the Stratified CIF curves is preserved, there is noticeable deviation across all risk strata compared to the empirical Aalen-Johansen estimates.

\begin{figure}[htbp]
    \centering
    \begin{subfigure}[b]{0.48\textwidth}
        \includegraphics[width=\textwidth]{images/pbc_results/SCC/Metrics/cif_calibration_cause_1_SCC_subcohort_0.5.png}
        \caption{Stratified CIF Calibration ($\alpha=0.5$)}
    \end{subfigure}
    \hfill
    \begin{subfigure}[b]{0.48\textwidth}
        \includegraphics[width=\textwidth]{images/pbc_results/SCC/Metrics/subject_curves_cause_1_cif_SCC_subcohort_0.5.png}
        \caption{Marginal Calibration ($\alpha=0.5$)}
    \end{subfigure}
    \caption{Calibration diagnostics for SCC with an optimal subcohort fraction ($\alpha=0.5$) on the PBC dataset.}
    \label{fig:pbc_scc_50_calibration}
\end{figure}

The consequence of these localized deviations is apparent in the Marginal Calibration plot (right panel of Figure \ref{fig:pbc_scc_50_calibration}). The mean model prediction systematically underestimates the population baseline, creating a persistent and expanding calibration gap. Unlike the synthetic dataset, which contained 15,000 observations, the PBC dataset is highly constrained. Therefore, discarding even 50\% of the background cohort removes critical covariate density required for accurate Inverse Probability Weighting (IPW), resulting in a dampened absolute risk profile.

This vulnerability becomes catastrophic under heavy compression. Figure \ref{fig:pbc_scc_05_calibration} illustrates the model's behavior when the background subcohort is reduced to just 5\% ($\alpha=0.05$). 

\begin{figure}[htbp]
    \centering
    \begin{subfigure}[b]{0.48\textwidth}
        \includegraphics[width=\textwidth]{images/pbc_results/SCC/Metrics/cif_calibration_cause_1_SCC_subcohort_0.05.png}
        \caption{Stratified CIF Calibration ($\alpha=0.05$)}
    \end{subfigure}
    \hfill
    \begin{subfigure}[b]{0.48\textwidth}
        \includegraphics[width=\textwidth]{images/pbc_results/SCC/Metrics/subject_curves_cause_1_cif_SCC_subcohort_0.05.png}
        \caption{Marginal Calibration ($\alpha=0.05$)}
    \end{subfigure}
    \caption{Calibration diagnostics for a heavily compressed SCC sampler ($\alpha=0.05$) demonstrating total IPW scaling failure.}
    \label{fig:pbc_scc_05_calibration}
\end{figure}

At $\alpha=0.05$, the IPW mechanism experiences a total structural collapse. The Stratified CIF curves (left panel) are effectively crushed to zero, failing entirely to track the empirical risk progression. Because the static subcohort is too sparse to represent the background hazard of the clinical population, the network optimizes against a distorted, near-zero baseline. This is confirmed by the Marginal Calibration plot (right panel), where the mean model prediction remains flat, creating a massive calibration gap that encompasses nearly the entire population incidence. 

This finding demonstrates that while static subcohorts offer computational stability in large datasets, they are mathematically fragile when applied to smaller clinical trials, as they risk destroying the absolute probability calibration of the subdistribution framework.

\subsubsection{Entropy Sampler}

The application of the dynamic Entropy sampler to the Primary Biliary Cirrhosis (PBC) dataset reveals a critical insight regarding hyperparameter tuning in the presence of real wolrd clinical noise. In the mathematically pure synthetic dataset, a highly greedy temperature scaling ($\tau=0.5$) acted as a powerful regularizer, filtering out redundant patients and improving risk recovery. However, applying this exact same greedy configuration to a small, noisy clinical dataset results in a severe calibration failure.

Figure \ref{fig:pbc_entropy_05_calibration} illustrates the collapse of the Entropy sampler at $\tau=0.5$. Because the PBC dataset is inherently noisy, forcing the sampler to hyper-fixate on the absolute highest-variance observations causes it to overfit to clinical outliers. This extreme sampling skew destroys the integrity of the Inverse Probability Weighting (IPW) mechanism. 

\begin{figure}[H]
    \centering
    \begin{subfigure}[b]{0.48\textwidth}
        \includegraphics[width=\textwidth]{images/pbc_results/Entropy/Metrics/cif_calibration_cause_1_Entropy_tau_0.5_compression_0.1.png}
        \caption{Stratified CIF Calibration ($\tau=0.5$)}
    \end{subfigure}
    \hfill
    \begin{subfigure}[b]{0.48\textwidth}
        \includegraphics[width=\textwidth]{images/pbc_results/Entropy/Metrics/subject_curves_cause_1_cif_Entropy_tau_0.5_compression_0.1.png}
        \caption{Marginal Calibration ($\tau=0.5$)}
    \end{subfigure}
    \caption{Calibration diagnostics for a highly greedy Entropy sampler ($\tau=0.5$) on the PBC dataset.}
    \label{fig:pbc_entropy_05_calibration}
\end{figure}

As seen in the left panel of Figure \ref{fig:pbc_entropy_05_calibration}, the predicted Stratified CIF curves systematically pull away and underestimate the true Aalen-Johansen incidence across all risk groups. Consequently, the Marginal Calibration plot (right panel) exhibits a massive, expanding calibration gap, as the mean model prediction entirely fails to capture the true baseline hazard of the population.

To correct this failure, the entropy of the sampling distribution must be artificially increased. By raising the temperature parameter to $\tau=5.0$, the sampling distribution becomes flatter and less greedy, forcing the model to retain a wider, more representative spread of the clinical cohort while still prioritizing information-dense observations. 

\begin{figure}[htbp]
    \centering
    \begin{subfigure}[b]{0.48\textwidth}
        \includegraphics[width=\textwidth]{images/pbc_results/Entropy/Metrics/cif_calibration_cause_1_Entropy_tau_5_compression_0.1.png}
        \caption{Stratified CIF Calibration ($\tau=5.0$)}
    \end{subfigure}
    \hfill
    \begin{subfigure}[b]{0.48\textwidth}
        \includegraphics[width=\textwidth]{images/pbc_results/Entropy/Metrics/subject_curves_cause_1_cif_Entropy_tau_5_compression_0.1.png}
        \caption{Marginal Calibration ($\tau=5.0$)}
    \end{subfigure}
    \caption{Calibration diagnostics for the optimal Entropy sampler ($\tau=5.0$), demonstrating complete restoration of absolute risk scaling.}
    \label{fig:pbc_entropy_5_calibration}
\end{figure}

As demonstrated in Figure \ref{fig:pbc_entropy_5_calibration}, this hyperparameter shift completely restores the model's absolute calibration. The Stratified CIF curves (left panel) closely track the empirical non-parametric estimates, and the Marginal Calibration Gap (right panel) narrows significantly, matching the structural integrity of the Full Cohort baseline. 

This behavior highlights the adaptive strength of the Entropy formulation: the temperature parameter $\tau$ acts as a tunable dial between computational compression and regularization. For large, clean datasets, it can be tuned down to aggressively compress the risk-set; for small, noisy clinical cohorts, it can be tuned up to preserve absolute probabilistic calibration without resorting to the brute-force computational overhead of a Full Cohort partial likelihood.

\newpage
\subsection{Real World Framingham Dataset}

To conclude the calibration analysis, the sampling architectures were evaluated on the Framingham heart study dataset. Unlike the highly distinct risk profiles of the synthetic data or the dense event rate of the PBC cohort, the Framingham dataset introduces a different biostatistical challenge: severe right-censoring. With a primary event incidence of only 15.6\%, the majority of the cohort is censored, creating an inherently sparse risk landscape over the long-term observational horizon.

\subsubsection{Full Cohort Baseline}

The Full Cohort configuration once again serves as the gold-standard baseline. By leveraging the entire available risk-set, it maximizes the amount of information extracted from the heavily censored data. 

As illustrated in Figure \ref{fig:fh_full_calibration}, the Full Cohort model handles the high censoring rate exceptionally well. The Stratified CIF Calibration (left panel) demonstrates that the model's predicted incidence correctly scales with empirical reality, even for the highest risk stratum, without experiencing probability inflation at the tail end of the survival curve. Furthermore, the Marginal Calibration diagnostic (right panel) shows near-perfect alignment between the mean model prediction and the population Aalen-Johansen baseline, with no visible calibration gap over the 3,500-day timeline.

\begin{figure}[htbp]
    \centering
    \begin{subfigure}[b]{0.48\textwidth}
        \includegraphics[width=\textwidth]{images/fh_results/FULL/Metrics/cif_calibration_cause_1.png}
        \caption{Stratified CIF Calibration}
    \end{subfigure}
    \hfill
    \begin{subfigure}[b]{0.48\textwidth}
        \includegraphics[width=\textwidth]{images/fh_results/FULL/Metrics/subject_curves_cause_1_cif.png}
        \caption{Marginal Calibration and Discrimination}
    \end{subfigure}
    \caption{Calibration diagnostics for the Full Cohort baseline on the heavily censored Framingham dataset.}
    \label{fig:fh_full_calibration}
\end{figure}

\newpage

\subsubsection{Nested Case-Control (NCC)}

The high censoring rate in the Framingham dataset exposes a unique interaction with the Nested Case-Control (NCC) sampler. In previous datasets, heavy risk-set compression ($m=5$) caused the IPW formulation to break down, resulting in a massive underestimation of absolute risk. However, as shown in Figure \ref{fig:fh_ncc_calibration}, the NCC architecture is remarkably resilient on this specific dataset, even at extreme sparsity.

\begin{figure}[H]
    \centering
    \begin{subfigure}[b]{0.48\textwidth}
        \includegraphics[width=\textwidth]{images/fh_results/NCC/Metrics/cif_calibration_cause_1_NCC_ncontrols_100.png}
        \caption{Stratified CIF Calibration ($m=100$)}
    \end{subfigure}
    \hfill
    \begin{subfigure}[b]{0.48\textwidth}
        \includegraphics[width=\textwidth]{images/fh_results/NCC/Metrics/subject_curves_cause_1_cif_NCC_ncontrols_100.png}
        \caption{Marginal Calibration ($m=100$)}
    \end{subfigure}
    
    \vskip 1\baselineskip
    
    \begin{subfigure}[b]{0.48\textwidth}
        \includegraphics[width=\textwidth]{images/fh_results/NCC/Metrics/cif_calibration_cause_1_NCC_ncontrols_5.png}
        \caption{Stratified CIF Calibration ($m=5$)}
    \end{subfigure}
    \hfill
    \begin{subfigure}[b]{0.48\textwidth}
        \includegraphics[width=\textwidth]{images/fh_results/NCC/Metrics/subject_curves_cause_1_cif_NCC_ncontrols_5.png}
        \caption{Marginal Calibration ($m=5$)}
    \end{subfigure}
    \caption{Calibration diagnostics for optimal and constrained NCC configurations on the Framingham dataset.}
    \label{fig:fh_ncc_calibration}
\end{figure}

Both the optimal configuration ($m=100$) and the heavily constrained configuration ($m=5$) successfully track the empirical Aalen-Johansen baseline without experiencing probability collapse. The Stratified CIF curves remain highly calibrated, and the Marginal Calibration Gap is negligible for both settings. 

This resilience is a direct artifact of the data structure. Because 84.4\% of the Framingham cohort is censored (meaning they do not experience an event), the background risk-set is highly uniform and dense with "survivors". When the NCC sampler draws $m=5$ random controls, those controls are almost certainly censored patients representing the true low-hazard baseline. Therefore, the uniform random sampling of NCC accurately approximates the population density without introducing the severe variance spikes that destroyed calibration in the denser PBC and Synthetic datasets.


\subsubsection{Static Case-Cohort (SCC)}

The application of the Static Case-Cohort (SCC) sampler to the Framingham dataset highlights a fundamental difference between static and dynamic sampling strategies under conditions of severe right-censoring. 

When a robust background cohort is maintained ($\alpha=0.5$), the SCC architecture performs optimally. As shown in the top row of Figure \ref{fig:fh_scc_calibration}, this configuration successfully replicates the Full Cohort baseline. The Stratified CIF curves track the empirical reality cleanly, and the Marginal Calibration plot shows near-perfect agreement between the model average and the population baseline.

\begin{figure}[htbp]
    \centering
    % Optimal Compression (alpha = 0.5)
    \begin{subfigure}[b]{0.48\textwidth}
        \includegraphics[width=\textwidth]{images/fh_results/SCC/Metrics/cif_calibration_cause_1_SCC_subcohort_0.5.png}
        \caption{Stratified CIF Calibration ($\alpha=0.5$)}
    \end{subfigure}
    \hfill
    \begin{subfigure}[b]{0.48\textwidth}
        \includegraphics[width=\textwidth]{images/fh_results/SCC/Metrics/subject_curves_cause_1_cif_SCC_subcohort_0.5.png}
        \caption{Marginal Calibration ($\alpha=0.5$)}
    \end{subfigure}
    
    \vskip 1\baselineskip
    
    % High Compression (alpha = 0.05)
    \begin{subfigure}[b]{0.48\textwidth}
        \includegraphics[width=\textwidth]{images/fh_results/SCC/Metrics/cif_calibration_cause_1_SCC_subcohort_0.05.png}
        \caption{Stratified CIF Calibration ($\alpha=0.05$)}
    \end{subfigure}
    \hfill
    \begin{subfigure}[b]{0.48\textwidth}
        \includegraphics[width=\textwidth]{images/fh_results/SCC/Metrics/subject_curves_cause_1_cif_SCC_subcohort_0.05.png}
        \caption{Marginal Calibration ($\alpha=0.05$)}
    \end{subfigure}
    \caption{Calibration diagnostics for SCC on the Framingham dataset, contrasting the optimal subcohort (top) with the failure of the highly compressed subcohort (bottom).}
    \label{fig:fh_scc_calibration}
\end{figure}

However, unlike the NCC sampler which proved highly resilient to extreme compression on this specific dataset the SCC sampler demonstrates a clear calibration breakdown when constrained to $\alpha=0.05$. In the bottom row of Figure \ref{fig:fh_scc_calibration}, the Stratified CIF curves for the higher risk strata begin to systematically underestimate the empirical Aalen-Johansen estimates. Consequently, the Marginal Calibration plot reveals a distinct and expanding calibration gap, as the mean model prediction detaches from the population baseline over time.

This divergence is mathematically illuminating. While the high censoring rate provided the NCC sampler with a dense pool of continuous "survivors" to sample from at each distinct event time, the SCC architecture is structurally rigid; it fixes its subcohort at $t=0$. In a long-horizon clinical study spanning nearly 10 years, drawing only a 5\% static sample at the baseline fails to provide sufficient covariate density to represent the evolving background hazard accurately. The IPW scaling becomes temporally skewed, proving that static sampling is more vulnerable to prolonged right-censoring than dynamic, time-matched sampling.

\subsubsection{Entropy Sampler}

The evaluation of the Entropy sampler on the heavily right-censored Framingham dataset exposes the complex interplay between dynamic risk-weighting and baseline hazard estimation. Because only 15.6\% of the Framingham cohort experiences the primary event, the background risk-set is overwhelmingly dominated by censored "survivors." 

When the temperature parameter is set to a highly greedy state ($\tau=0.5$), the Entropy algorithm aggressively prioritizes the rare, high-variance patients experiencing events, severely down-weighting the massive pool of censored observations. In a highly censored dataset, this over concentration proves detrimental to absolute calibration. 

\begin{figure}[htbp]
    \centering
    \begin{subfigure}[b]{0.48\textwidth}
        \includegraphics[width=\textwidth]{images/fh_results/Entropy/Metrics/cif_calibration_cause_1_Entropy_tau_0.5_compression_0.1.png}
        \caption{Stratified CIF Calibration ($\tau=0.5$)}
    \end{subfigure}
    \hfill
    \begin{subfigure}[b]{0.48\textwidth}
        \includegraphics[width=\textwidth]{images/fh_results/Entropy/Metrics/subject_curves_cause_1_cif_Entropy_tau_0.5_compression_0.1.png}
        \caption{Marginal Calibration ($\tau=0.5$)}
    \end{subfigure}
    \caption{Calibration diagnostics for a greedy Entropy sampler ($\tau=0.5$) on the Framingham dataset, showing systemic over- and underestimation of risk.}
    \label{fig:fh_entropy_05_calibration}
\end{figure}

As illustrated in Figure \ref{fig:fh_entropy_05_calibration}, the $\tau=0.5$ configuration distorts the Inverse Probability Weighting (IPW) scaling. In the left panel, the model noticeably overestimates the true incidence for the high-risk stratum (solid green line exceeding the dashed empirical line). Concurrently, the Marginal Calibration diagnostic (right panel) reveals a distinct calibration gap, with the mean model prediction systematically underestimating the overall population baseline. By starving the neural network of the censored background population, the model fails to accurately anchor its absolute probability estimates.

To rectify this, the temperature parameter must be increased to soften the sampling distribution. Setting $\tau=5.0$ forces the Entropy sampler to behave less greedily, ensuring that while high-information events are still prioritized, a statistically viable volume of censored background patients is retained in every mini-batch. 

\begin{figure}[htbp]
    \centering
    \begin{subfigure}[b]{0.48\textwidth}
        \includegraphics[width=\textwidth]{images/fh_results/Entropy/Metrics/cif_calibration_cause_1_Entropy_tau_5_compression_0.1.png}
        \caption{Stratified CIF Calibration ($\tau=5.0$)}
    \end{subfigure}
    \hfill
    \begin{subfigure}[b]{0.48\textwidth}
        \includegraphics[width=\textwidth]{images/fh_results/Entropy/Metrics/subject_curves_cause_1_cif_Entropy_tau_5_compression_0.1.png}
        \caption{Marginal Calibration ($\tau=5.0$)}
    \end{subfigure}
    \caption{Calibration diagnostics for the optimally tuned Entropy sampler ($\tau=5.0$) demonstrating the closure of the Marginal Calibration Gap.}
    \label{fig:fh_entropy_5_calibration}
\end{figure}

Figure \ref{fig:fh_entropy_5_calibration} demonstrates the complete restoration of absolute calibration at $\tau=5.0$. The Stratified CIF curves now align perfectly with the empirical Aalen-Johansen estimates across all strata, eliminating the high-risk overestimation observed at $\tau=0.5$. Furthermore, the Marginal Calibration Gap is entirely closed, matching the structural integrity of the Full Cohort baseline.

This final evaluation solidifies the overarching operational framework for the CRISPNAM architecture. Across both small, noisy clinical trials (PBC) and heavily censored longitudinal studies (Framingham), highly greedy entropy sampling ($\tau=0.5$) corrupts IPW calibration. However, by treating the temperature parameter $\tau$ as an environmental hyperparameter tuning it upward ($\tau \in [1.0, 5.0]$) to flatten the sampling distribution in the presence of severe data sparsity or noise the Entropy sampler successfully matches Full Cohort predictive power and absolute calibration at a fraction of the computational cost.

\section{Feature Contribution and Interpretability}

Beyond absolute predictive calibration, the synthetic dataset provides a controlled environment to evaluate the interpretability of the CRISPNAM architecture. Because the underlying data-generating parameters ($\beta$) are known, the isolated shape functions $f_i(x_i)$ learned by the Neural Additive Model (NAM) backend can be visually mapped directly against the true biostatistical feature effects.

Figure \ref{fig:feature_contributions_cause_1} illustrates this component-wise recovery for the primary event of interest (Cause 1) using the Full Cohort baseline. The solid blue lines represent the empirical shape functions learned entirely from the data, while the dashed red lines denote the true underlying linear hazard contributions.

\begin{figure}[H]
    \centering
    \includegraphics[width=0.85\textwidth]{images/Synth_results/FULL/feature_contributions_Event_1_cause_1.png}
    \caption{Isolated FeatureNet contributions for Cause 1 (Synthetic Dataset, Full Cohort).}
    \label{fig:feature_contributions_cause_1}
\end{figure}

\subsection{Signal Recovery and Noise Isolation}

The diagnostic plots demonstrate exceptional feature disentanglement. For strong true-signal features (e.g., $x_1, x_2, x_{11}, x_{12}$), the learned nonlinear shape functions align closely with the true linear trajectories across the core feature domains, accurately capturing the magnitude and directionality of the hazard contributions. 

Furthermore, the plots visually validate the architectural use of the Sigmoid Linear Unit (SiLU) activation function. Unlike bounded activations (such as $\tanh$), which artificially saturate and flatten out at the extremes of a feature's distribution, the unconstrained nature of the SiLU activation allows the learned shape functions to maintain smooth, continuous trajectories at the outer bounds of the covariate space (most clearly visible at the extremes of $x_1$ and $x_2$).

Crucially, the additive framework successfully isolates observational noise. For features where the true generating parameter is zero (e.g., $x_3, x_4, x_9, x_{10}$), the network learns minor non-linear fluctuations rather than a perfectly flat line. This is an expected artifact of local overfitting within a finite synthetic sample. However, because each feature is processed independently, this spurious noise is strictly quarantined. The noise captured in $f_3(x_3)$ cannot entangle with or corrupt the pristine signal recovered in $f_1(x_1)$. This confirms that the CRISPNAM architecture functions as a transparent, white-box survival model, clearly revealing the exact drivers of patient-level risk.


\subsection{Impact of Risk-Set Sampling on Feature Recovery}

The introduction of risk-set sampling alters the gradient landscape during backpropagation. While the previous sections established how these samplers affect absolute probabilistic calibration, it is equally critical to understand how fractional data exposure impacts the isolated FeatureNets' ability to learn the true underlying biostatistical shapes. 

By analyzing the component-wise feature recovery across the samplers, the previously observed macro-level metrics (C-index, Pearson, and Spearman correlations) can be visually explained at the network's foundational level.

\subsubsection{Dynamic Entropy Sampler}

The optimal Entropy sampler ($\tau=0.5$) provides the most mathematically profound finding of this architecture. Rather than degrading feature recovery due to data omission, the intelligent targeting of high-variance subjects acts as a powerful data-level regularizer. 

As illustrated in Figure \ref{fig:feature_recovery_entropy}, the optimal Entropy configuration actually reconstructs the true generating parameters ($\beta$) with higher fidelity than the Full Cohort baseline. The learned shape functions for the active features (such as $x_1$, $x_2$, $x_7$, and $x_{12}$) exhibit smoother, more precise trajectories that tightly hug the dashed ground-truth lines. This visual precision directly correlates with the superior macro-metrics achieved by this configuration: a peak Pearson correlation of 0.866, a Spearman rank of 0.857, and an optimal C-index of 0.851. By aggressively down-weighting dense clusters of redundant, low-information patients, the Entropy sampler prevents the FeatureNets from overfitting to local population noise, allowing them to isolate a cleaner hazard signal.

\begin{figure}[htbp]
    \centering
    \begin{subfigure}[b]{0.48\textwidth}
        \includegraphics[width=\textwidth]{images/Synth_results/ENTROPY/feature_contributions_Event_1_cause_1_Entropy_tau_0.5_compression_0.1.png}
        \caption{Optimal Entropy ($\tau=0.5$)}
    \end{subfigure}
    \hfill
    \begin{subfigure}[b]{0.48\textwidth}
        \includegraphics[width=\textwidth]{images/Synth_results/ENTROPY/feature_contributions_Event_1_cause_1_Entropy_tau_5_compression_0.1.png}
        \caption{Relaxed Entropy ($\tau=5.0$)}
    \end{subfigure}
    \caption{Feature recovery for the Entropy sampler comparing optimal greedy behavior to relaxed uniform sampling.}
    \label{fig:feature_recovery_entropy}
\end{figure}

Conversely, as the temperature parameter is artificially relaxed ($\tau=1.0$ moving to $\tau=5.0$), the sampling distribution becomes flatter and more uniform. Consequently, the feature recovery gracefully degrades back toward the Full Cohort baseline. At $\tau=5.0$ (right panel of Figure \ref{fig:feature_recovery_entropy}), minor non-linear distortions begin to reappear in the zero-signal features, and the Pearson correlation drops to 0.839 (with a C-index of 0.848), confirming that uniform random inclusion reintroduces the observational noise that the greedy configuration successfully filtered out.

\subsubsection{Nested Case-Control (NCC)}

The Nested Case-Control architecture's feature recovery is strictly bounded by the size of its control pool ($m$). When operating with a large quota ($m=100$), the NCC model behaves almost identically to the Full Cohort. It achieves a high Pearson correlation of 0.856 and a C-index of 0.849, and its corresponding feature plots successfully capture the correct directionality and magnitude of the covariates.

\begin{figure}[htbp]
    \centering
    \begin{subfigure}[b]{0.48\textwidth}
        \includegraphics[width=\textwidth]{images/Synth_results/NCC/feature_contributions_Event_1_cause_1_NCC_ncontrols_100.png}
        \caption{Optimal NCC ($m=100$)}
    \end{subfigure}
    \hfill
    \begin{subfigure}[b]{0.48\textwidth}
        \includegraphics[width=\textwidth]{images/Synth_results/NCC/feature_contributions_Event_1_cause_1_NCC_ncontrols_5.png}
        \caption{Constrained NCC ($m=5$)}
    \end{subfigure}
    \caption{Feature recovery for the NCC sampler demonstrating signal dampening under extreme constraint.}
    \label{fig:feature_recovery_ncc}
\end{figure}

However, extreme compression fundamentally starves the backpropagation gradients. As shown in the right panel of Figure \ref{fig:feature_recovery_ncc}, restricting the network to $m=5$ controls induces visible distortion. While the network correctly identifies the positive or negative trajectory of the active features (preserving a baseline C-index of 0.835), it fails to learn the true magnitude of the slopes. The curves for $x_1$ and $x_2$ pull away from the true $\beta$ trajectories, and the noise within the zero-signal features ($x_3$, $x_{10}$) becomes highly amplified. This geometric distortion explains the severe drop in the linear Pearson correlation (0.806) and the Spearman rank (0.795). 

\subsubsection{Static Case-Cohort (SCC)}

The Static Case-Cohort (SCC) sampler demonstrates a similar threshold-based degradation, though its fixed background cohort provides slightly more gradient stability than the highly stochastic constrained NCC sampler. 

\begin{figure}[htbp]
    \centering
    \begin{subfigure}[b]{0.48\textwidth}
        \includegraphics[width=\textwidth]{images/Synth_results/SCC/feature_contributions_Event_1_cause_1_SCC_subcohort_0.5.png}
        \caption{Optimal SCC ($\alpha=0.5$)}
    \end{subfigure}
    \hfill
    \begin{subfigure}[b]{0.48\textwidth}
        \includegraphics[width=\textwidth]{images/Synth_results/SCC/feature_contributions_Event_1_cause_1_SCC_subcohort_0.05.png}
        \caption{Constrained SCC ($\alpha=0.05$)}
    \end{subfigure}
    \caption{Feature recovery for the SCC sampler comparing optimal subcohort retention to extreme compression.}
    \label{fig:feature_recovery_scc}
\end{figure}

At an optimal subcohort fraction ($\alpha=0.5$), the feature plots (left panel of Figure \ref{fig:feature_recovery_scc}) show strong alignment with the ground truth, matching the Full Cohort baseline with a Pearson score of 0.849 and a C-index of 0.849. However, when the static subcohort is decimated to just 5\% of the population ($\alpha=0.05$), the FeatureNets lose the necessary density to map the extremes of the covariate space. Consequently, the learned shape functions flatten out and dampen, leading to a measurable drop in linear recovery (Pearson 0.835, Spearman 0.825, C-index 0.842). 

Ultimately, this component-level visual analysis confirms that only the dynamic Entropy sampler is capable of compressing the risk-set footprint without mathematically dampening or distorting the underlying biostatistical shape functions.


\section{Discussion}

The primary objective of this research was to resolve the computational bottleneck of deep survival architectures under competing risks without sacrificing absolute probability calibration or feature interpretability. By systematically evaluating the CRISPNAM architecture across pure mathematical simulations (Synthetic), standard clinical trials (PBC), and heavily right-censored longitudinal studies (Framingham), this study provides a comprehensive blueprint for optimizing deep survival networks. 

The results generated validate the architecture's efficacy as both a highly interpretative survival model and a computationally robust framework for modeling competing risks.

\subsection{Interpretation of Findings}

The overarching takeaway from the empirical evaluations is that dynamic, information-aware risk-set sampling strictly outperforms static or uniform sampling methods in deep survival analysis. The traditional algorithms Nested Case-Control (NCC) and Static Case-Cohort (SCC) exhibited severe structural vulnerabilities. When the active risk-sets were aggressively compressed to maximize computational speed, both NCC and SCC suffered from catastrophic calibration failures. Because deep neural networks optimize via mini-batch gradient descent, starving the network of the background population skewed the Inverse Probability Weighting (IPW) mechanisms, causing the models to systematically under- or overestimate the absolute Cumulative Incidence Functions (CIF).

Conversely, the Entropy sampler demonstrated that the active risk-set can be heavily compressed without sacrificing probabilistic calibration or feature recovery, provided the sampling distribution targets high-variance, information-dense subjects. On the synthetic dataset, the optimal Entropy configuration ($\tau=0.5$) actually acted as a data-level regularizer, recovering the true biostatistical generating parameters ($\beta$) with a higher Pearson correlation (0.866) and C-index (0.851) than the Full Cohort baseline. 

Furthermore, the evaluation revealed a critical hyperparameter dynamic: the temperature scaling ($\tau$) must be actively tuned to the environment's noise profile. In mathematically clean datasets, a greedy approach ($\tau=0.5$) isolates pristine signals. However, in noisy clinical datasets (PBC) or heavily censored cohorts (Framingham), relaxing the temperature ($\tau \in [1.0, 5.0]$) becomes necessary to prevent the network from overfitting to clinical outliers. 

Crucially, the Neural Additive Model (NAM) backend upgraded with SiLU activations to prevent boundary saturation ensured that these computational gains did not come at the cost of interpretability. The architecture functioned as a transparent white-box model, allowing the direct visualization of isolated, patient-level hazard drivers without cross-feature entanglement.

\subsection{Comparative Validation Against Existing Literature}

To ensure the structural integrity of the underlying Fine-Gray loss formulation prior to sampling, the Full Cohort baselines were benchmarked against contemporary deep survival literature. Specifically, the model's performance was compared to the benchmark results established in the \textit{Interpretable Fine-Gray Deep Survival Model for Competing Risks} \citep{InterpretableFineGray2025}. 

To facilitate a direct, one-to-one comparison with the referenced study, the Time-Dependent Concordance Index (TD-CI) and Brier scores for the clinical cohorts were explicitly extracted at the 25th ($q_{0.25}$), 50th ($q_{0.50}$), and 75th ($q_{0.75}$) percentiles of the temporal distribution, alongside the overall pooled metrics.

For the primary event in the PBC dataset, the referenced literature reports a median TD-CI ($q_{0.50}$) of $0.813 \pm 0.016$ and a median Brier Score of $0.249 \pm 0.089$ for the CRISPNAM-FG architecture. The Full Cohort model in this study achieved a $q_{0.50}$ TD-CI of $0.663$ (peaking at $0.688$ for the $q_{0.75}$ horizon) and a $q_{0.50}$ Brier Score of $0.180$. While the discriminative rank-ordering is lower than the benchmark—a direct reflection of our rigorous patient-grouped cross-validation that intentionally eliminates optimistic within-subject data leakage the vastly superior (lower) Brier score achieved in this study confirms that the absolute probabilistic calibration of the architecture remains exceptionally sound.

On the Framingham Heart Study (FHS) dataset, the referenced literature reports a median TD-CI of $0.755 \pm 0.008$ and a median Brier Score of $0.202 \pm 0.141$. The Full Cohort model in this study recorded a $q_{0.50}$ TD-CI of $0.607$ and a $q_{0.50}$ Brier Score of $0.670$. This deviation from the reference benchmarks is a direct and expected result of methodological divergence in dataset preparation. The FHS dataset utilized in this study underwent a distinctly different feature engineering pipeline and was modified to forcefully induce a competing risk scenario. These structural alterations naturally shift the probabilistic baselines and covariate topographies. However, achieving a stable pooled C-index of $0.645$ using the full temporal evaluation confirms that the model successfully isolates the underlying clinical hazard despite the induced competing events and the extreme sparsity of actual events (84.4\% censored).

Finally, on the robust synthetic simulation, the Full Cohort achieved a pooled C-index of $0.850$ and an Brier Score of $0.106$. This proves that absent severe clinical noise, the mathematical foundations of the network and the Fine-Gray objective function operate at capacity, successfully establishing the stable theoretical baseline required to conduct the advanced riskset sampling experiments.

\subsection{Computational Efficiency and Vectorization}

The profiling of the sampling models revealed a critical paradox in the application of traditional risk-set sampling (NCC and SCC) within modern deep learning frameworks. Historically, nested case-control and case-cohort sampling were developed to alleviate the computational burden of standard, single-event Cox models running sequentially on CPU architectures. 

However, modern deep neural networks rely fundamentally on highly parallelized matrix multiplications optimized for Graphical Processing Units (GPUs). As observed during the evaluation of the small Framingham cohort, the Full Cohort model actually trained faster than the highly constrained NCC ($m=5$) model. This occurs because calculating the partial likelihood for the full cohort requires continuous, unbroken tensor operations that easily fit into VRAM, maximizing GPU throughput. In contrast, standard NCC and SCC methods require dynamic subsampling loops to construct the risk-sets at each unique event time. This introduces severe CPU-to-GPU data transfer bottlenecks, breaking the inherent vectorization of the computational framework. The operational overhead of naive sampling outpaces the speed gained by processing smaller matrices.

This hardware dynamic mathematically explains why the Entropy sampler proves vastly superior. By calculating the sampling probabilities upfront and intelligently reducing the data footprint based on information density rather than uniform algorithmic loops, the Entropy sampler maintains the structural integrity of the active tensors. On the large synthetic dataset where the full cohort tensor eventually outgrows efficient VRAM handling the Entropy sampler cut execution time by over 50\% without breaking GPU vectorization. This demonstrates that for deep survival architectures, intelligent, regularized data compression is computationally superior to traditional uniform control matching.

\subsection{Theoretical and Practical Implications}

The findings of this study carry significant implications for the intersection of biostatistics and deep learning. Theoretically, it establishes that deep survival architectures should not default to traditional epidemiological sampling methods (like NCC or SCC) to solve computational bottlenecks. Because neural networks require robust feature variance to map non-linear functions effectively, uniform or static sampling actively harms the gradient landscape. Instead, sampling in deep survival analysis must be viewed as an information theory problem, where data compression is guided by variance and entropy.

Practically, the CRISPNAM architecture with Entropy sampling provides a highly scalable solution for modern healthcare analytics. As Electronic Health Records (EHR) generate massive, highly dimensional datasets, calculating the full partial likelihood becomes physically intractable on standard clinical hardware. The proposed architecture allows data science and biostatistics teams to train highly accurate competing risk models on massive patient populations using standard computational resources. Furthermore, the NAM backend ensures that these models can be deployed in clinical settings without violating the need for medical transparency, as physicians can visually verify the exact biomarkers driving a patient's risk score.


\chapter{Conclusion and Future Work}

\section{Summary of Findings}

This dissertation addressed the critical computational bottlenecks inherent in training deep survival architectures for competing risks. By introducing and rigorously evaluating the CRISPNAM architecture alongside dynamic risk-set sampling mechanisms, this research demonstrated that computational efficiency does not have to come at the expense of absolute probabilistic calibration or clinical interpretability. 

The empirical evaluations across pure synthetic simulations and real wolrd clinical datasets (PBC, Framingham) conclusively showed that traditional sampling techniques Nested Case-Control (NCC) and Static Case-Cohort (SCC) exhibit severe structural vulnerabilities in deep learning contexts. While they successfully reduce the data footprint, they disrupt the continuous tensor operations required for GPU vectorization and routinely destroy Inverse Probability Weighting (IPW) calibration under conditions of extreme risk-set sparsity. 

Conversely, the dynamic Entropy sampler proved vastly superior. By calculating sampling probabilities upfront and prioritizing information-dense, high-variance subjects, the Entropy mechanism effectively acts as a robust data-level regularizer. The evaluations revealed that when the temperature parameter ($\tau$) is actively tuned to the dataset's intrinsic noise profile using highly greedy sampling ($\tau \le 0.5$) for clean datasets and softer distributions ($\tau \ge 1.0$) for noisy or heavily censored clinical cohorts the model thrives. Under optimal configurations, the Entropy sampler matched or exceeded the Full Cohort baseline's predictive performance (C-index, Brier Score) and ground-truth feature recovery, while reducing computational training time by over 50\%. 

Crucially, the integration of a Neural Additive Model (NAM) backend enhanced with SiLU activations to prevent boundary saturation ensured that the architecture remained fully transparent. The model successfully disentangled complex covariate effects, allowing for the direct visualization of individual, patient-level hazard drivers without the "black-box" entanglement characteristic of standard deep survival networks.

\section{Limitations and Future Directions}

Despite the robust findings, several limitations inherent to the study's scope and methodology present immediate, highly promising avenues for future research.

First, the reliance on standard IPW to correct sampling bias proved structurally fragile under extreme risk-set compression. As observed in the SCC and heavily constrained NCC evaluations, excessive subsampling caused the IPW denominators to approach zero, failing to represent the background hazard and resulting in immediate probability collapse. Future work should investigate the integration of stabilized IPW and Inverse Probability of Censoring Weighting (IPCW) mechanisms. Furthermore, coupling these stabilized weighting techniques with the Entropy sampling framework and applying them to other state-of-the-art deep survival architectures (e.g., DeepHit, DeepSurv) would verify whether these computational and predictive gains are universally transferable across different neural network topologies.

Second, while this study explicitly focused on the Fine-Gray subdistribution hazard for competing risks, the core mechanism of variance-driven risk-set compression is theoretically agnostic to the event type. Future research should evaluate whether this intelligent Entropy sampling technique yields the same computational acceleration and regularization advantages in standard, single-event survival analysis (e.g., traditional Cox Proportional Hazards models).

Third, the temperature parameter ($\tau$) in the Entropy sampler currently functions as a static environmental hyperparameter. Finding the optimal $\tau$ requires a grid search prior to training, adding overhead to the model development lifecycle. A critical future direction involves developing an adaptive, learnable temperature scaling mechanism. Designing $\tau$ to decay dynamically across optimization epochs starting highly uniform ($\tau > 5.0$) to securely anchor the IPW baseline hazard, and cooling down to a greedy state ($\tau \le 0.5$) in later epochs to fine-tune signal recovery would fully automate the architecture's adaptation to dataset noise.

Finally, the upfront computational cost of calculating the entropy distribution across the risk-sets scales alongside dataset size. While this proved highly efficient overall in the evaluated cohorts, expanding the sampling methodology to Deep Multi-State Models presents a rigorous future challenge. Multi-state modeling provides a more comprehensive overview of disease progression than competing risks alone but suffers from even more severe computational scaling limits. Applying the dynamic Entropy sampler to large-scale multi-state transition matrices represents the next logical frontier for scaling complex longitudinal biostatistical modeling in modern healthcare analytics.

\bibliography{references}\addcontentsline{toc}{chapter}{References}
\end{document}